\subsection{The coefficient matrix on the selected evaluation grid}
\label{subsec:fitted-evaluation}

The square ratios account for the scalar between tuples. A matrix of algebraic coefficient evaluations is still needed to identify their row images over the coefficient field.

\begin{paramdefinition}[Indices and fields of a fitted grid]\label{hyp:fitted-grid-data}
Fix the weight parameters, the grid cutoff \(C_F\), the base tuple
\(t_0\) and the frame of the selected row. For each finite coefficient
place \(w\) of the larger field \(K\), the grid uses
\[
     w_0=w|_{K_0},\qquad t(w)=\text{the tuple assigned to }w_0 ,
       \qquad K_0\subset K .
\]
Use the Weil model and smooth test vector of this row for its coefficient
integrals. Proposition~\ref{prop:fitted-grid-field-fit} constructs
a common algebraic field for the finitely many values required below.
\end{paramdefinition}

\subsubsection{A common field and its test models}

The local tests and their conjugates involve several coefficient models. Their fields must be fitted simultaneously before an entry can be interpreted in the chosen rational basis.

\begin{proposition}[Algebraic fields for the occurring grids]\label{prop:fitted-grid-field-fit}
Use Definition~\ref{hyp:fitted-grid-data}. Fix the weight
and lower-level data before choosing the coefficient fields
and the finite list of allowed tuple characters.
There are finitely many algebraic inclusions which can be imposed
at the field-choice stage so that the character and prescribed
root values, jet entries, terminal pairings and critical-value
ratios all belong to the larger field \(K\). Let \(S_0\) consist of the values of the column characters and the stipulated base roots of unity, and let \(S_1\) consist of the values listed above for \(K\). At the chosen embeddings these finite requirements are
\[
 S_0\subset K_0\subset K,\qquad S_1\subset K.
\]
 Transport preserves these values on each prescribed grid
with its prescribed tuple and levels.
\end{proposition}

\begin{proof}
\proofstep{Count the finite labels.}

Let \(\mathcal X\) be the finite list of allowed tuple
characters. For each \(\chi\in\mathcal X\), its extra
level is a subset of the finite support \(S(\chi)\) of
ramified places, and \(S(\chi)\) lies in that level and
its conjugate. Hence the first grid label belongs to
the subsets of a set of at most
\(\sum_{\chi\in\mathcal X}2^{|S(\chi)|}\) character-level pairs.

For the second label, the split test place divides the fixed
integer \(q_0\), so there are finitely many pairs of that place
and a place of \(E\) above it. At fixed infinity type any two
admitted column characters have quotient trivial at infinity
and of conductor dividing the level \(\mathfrak m_{\rm col}\)
with exponent two at the allowed finite places and zero
elsewhere. If one such character \(\eta_0\) exists, the
others belong to
\[
       \eta_0\,\operatorname{Cl}_{\mathfrak m_{\rm col}}(E)^\vee.
\]
The ray class character group is finite. Applying this count
at each column of the chain gives finitely many second labels.

\proofstep{Choose fields of definition for the normalized values.}

For each occurring pair of labels choose a representative grid
and real test. Work over the number field generated by the
following finite lists.

\begin{enumerate}
\item Take the algebraic column-character values, the required
primitive roots of unity and positive square roots of residue
cardinalities. Algebraic Hecke character values come from
their principal-idele formula and representatives of a finite
ray class quotient. At a place of conductor exponent \(a>0\),
the character on local units factors through

\[
  (\mathcal O_v/\mathfrak p_v^a)^\times.
\]
 It is trivial on
units at the other finite places. There are thus only finitely
many local unit values to adjoin. Normalization to the
prescribed algebraic infinity exponents is included in this choice.
\item Expand each coset term at its prescribed depth. The data
and tuple set are finite, and each coset term is expanded using
a fixed finite index set at that depth. Write \(D_t\) for the allowed indices at depth \(t\). The support of every such expansion lies in the finite union
\[
  \bigcup_{t\ {\rm used}}D_t,\qquad \#\bigcup_{t\ {\rm used}}D_t\le\sum_{t\ {\rm used}}\# D_t .
\]

For each occurring definite class and both chains take the
entries of the normalized jet in the coefficient bases. The
entries are algebraic by the CM normalization of the differential
and the jet comparison at the same finite level
(Section~\ref{subsec:selected-kernel-jets}).
\item Take the finite values of both terminal triangles and
the critical-value ratios. Each bottom lift is spanned by class
sums supported on finitely many representatives and by the
monomials at its fixed weight. The untwisted and twisted
normalized sums are algebraic by the level and jet comparisons.
The two bottom spaces can have different generating lists.
Proposition~\ref{prop:fitted-square-ratio-identity} gives the
algebraic ratios. Toric transport with the common period gives
the values of the two deletion terms. The normalized
values in each list are attached to the representative's
level and infinity exponents.
\end{enumerate}

\proofstep{Transport the values through the level identifications.}

For each occurring label pair \(\lambda\), use the finite
list \(A_\lambda\) just constructed on its representative.
Set 
\[
  K_\lambda=\mathbb Q(A_\lambda)
\]
 and choose a rational
basis \(\mathcal B_\lambda\). The entire request is finite:
\[
 \bigcup_\lambda\mathcal B_\lambda\ \subset\ \iota(K)
       \quad\Longrightarrow\quad
 K_\lambda\subset\iota(K)\quad\text{for every occurring }\lambda .
\]
Indeed each \(K_\lambda\) is generated by algebraic numbers
and only finitely many label pairs occur.

The comparison of equal labels identifies the actual tuple
and its representative, including its character, extra level
and local models. In each finite coset expansion write a
representative term as \(a\,T_\xi \Phi\), with \(T_\xi\)
the prescribed tensor translate. Transport carries \(T_\xi\)
to the translate at the identified level and multiplies its
coefficient by the specified additive phase, a root already
in \(K_\lambda\). Its depth and support are unchanged
in the jet comparison.

For a term in this index set with original coefficient \(a_\xi\) this transport changes its valuation only by that of the additive phase \(\zeta_\xi\), so
\[
  v(\zeta_\xi a_\xi)=v(a_\xi),\qquad v(\zeta_\xi)=0 .
\]
 This identifies each normalized jet
entry and class sum with the transported expression indexed by the same finite set. The two terminal triangles use their
respective transported bases. Thus these field requirements
apply before choosing an evaluation point, and include the
deletion values on either side. The smaller field contains
the column-character values and at each base place the roots
of order \((\mathrm Nv)^{h_v}\) specified by the depth \(h_v\).
The normalized matrix entries use \(K\).

\label{cor:fitted-evaluation-1}

\end{proof}

\subsubsection{Construction of the coefficient matrix}

With the field inclusions fixed, evaluate the finite class sums in the prescribed input basis. This identifies the matrix entries used at both finite and complex places.

\begin{proposition}[Algebraicity in the selected input basis]\label{prop:fitted-grid-coefficient-matrix}
With Proposition~\ref{prop:fitted-grid-field-fit}, the normalized
integrals for the prescribed coefficient tests and the selected
rational input basis lie in \(\sigma_0(K)\). Writing \(I_{ab}(t)\) for the integral in the \(a\)-th test on the \(b\)-th basis vector, the coefficient matrix is characterized by
\[
 V(t)_{ab}\in K,\qquad \sigma_0(V(t)_{ab})=I_{ab}(t).
\]
The coefficient identity extends to every rational tuple vector with the prescribed conjugate embedding in the input realization.
\end{proposition}

\begin{proof}
\proofstep{Use the selected rational tuple basis.}

Let \(u_a\) be the finitely many coefficient tests, indexed
by the datum and definite-class blocks and by a basis in each
block. Use the input column basis at the actual selected level
to form \(e_{t,j}\), the corresponding vectors of the rational
tuple space. For each tuple in use set
\[
    V(t)_{aj}=\sigma_0^{-1}
        \bigl(\mathcal V_t(u_a,e_{t,j})\bigr),
\]
where \(\mathcal V_t\) is the normalized coefficient integral in
the field coordinates and \(\sigma_0:K\hookrightarrow\C\) is the
specified embedding.

\proofstep{Read both chains at the common terminal class.}
Unfold the integral by the matched local factors. The
finite coset and class expansion is a sum of products of the
local coefficients, the indexed jet entries, the rational
input-form and test coordinates, the rational terminal pairing
and the external scalar. Its factors are in the coefficient fields just constructed. The rational inputs use the Lie frames of
the actual class points. The scalar is a product of integer
powers of the chosen square roots times rational factors.
The finite expansion uses both chains at the same terminal
class, with the pairing after their coefficient maps.

The terminal pairing therefore equals the embedded sum of field products.
In particular a finite sum of coefficient products commutes with the chosen embedding \(\iota\)
\[
  \iota\!\left(\sum_c a_c b_c\right)=\sum_c \iota(a_c)\iota(b_c)\qquad(a_c,b_c\in\operatorname{dom}\iota).
\]
 The
coefficient identity is linear in the embedded rational input
coordinates, with the prescribed conjugate embedding built into
the realization of that input. The calculation on the basis
therefore proves the coefficient identity for every vector in
the rational tuple space. Only \(t_0\) and the tuples assigned to
finite places enter the local and complex estimates.
\label{cor:fitted-evaluation-2}

\end{proof}

\subsubsection{Separation at the base}

Algebraicity of the entries alone does not control the rank of the resulting row. The base tests must separate the quotient on which the period is defined.

\begin{paramdefinition}[Coefficient functionals on a specified row]\label{hyp:fitted-pointwise-nonvanishing}
At a row scale \(x\) with error parameter \(\varepsilon\) and cutoff
\(C_F\ge6510\), use the grid of Definition~\ref{hyp:fitted-grid-data}.
Let \(w\ne0\) be the chosen decomposable smooth vector and \(V\) the
base multiplicity space. For a column \(\Psi\) of the principal level,
kernel datum \(\Phi\), real vector \(\rho\) and coefficient covector \(\xi\),
denote the coefficient integral against the row \(w,v\) by
\[
         v\longmapsto I_v(\rho,\xi;\Psi,\Phi),\qquad v\in V.
\]
Lemma~\ref{lem:fitted-completed-nonvanishing} proves that for
\(x\ge x_0(\varepsilon,C_F)\) these functionals separate the vectors
of \(V\); the test may depend on the vector.
\end{paramdefinition}

The auxiliary grid in Definition~\ref{hyp:fitted-pointwise-nonvanishing}
contains a row \(Y_i\) and a companion \(C_i\) at each stage \(0\le i\le m\).
If \(r=\dim_E Y_0\), their ranks are
\[
 \dim_E Y_i=r+i(r+6),\qquad \dim_E C_i=r+1+i(r+5).
\]
The bottom companion is \(C_0=Y_0\perp\langle1\rangle\); it has
the same negative index as \(Y_0\), namely \(m/2\) at the selected
real places. The last companion \(C_m\) is definite. Call a bottom
form admissible if it is obtained by descending from an allowed column
on \(U(C_m)\) along these companion spaces. The bottom period pairs
that form on \(U(C_0)\) with the given row on \(U(Y_0)\). Use the
measures and character twists of Lemma~\ref{hyp:joint-period-towers}.

The depth in the pointwise question matters already at the bottom of
the tower. One can prove nonvanishing in a family of deeper principal
levels by an approximation on the unitary group. To see precisely where
a bound for the depth would enter, fix the Hermitian spaces first.

\begin{auxlemma}[A bottom period at some finite depth]\label{lem:bottom-unbounded-depth}
For the row space \(Y\) put
\[
 C=Y\perp\langle1\rangle,\qquad H=U(Y)\subset G=U(C),
 \qquad n=\dim_E C.
\]
At a split finite place \(v\) allow the level to vary. Fix compact
open levels satisfying
\[
 J_H=J_H^vJ_{H,v},\qquad K^v\subset G(\mathbb A_{F,f}^v),\qquad
 K^v\cap H(\mathbb A_{F,f}^v)\subset J_H^v .
\]
At \(v\) fix an integral frame and its principal subgroups \(K_v(d)\).
Suppose \(\sigma\) and \(\pi\) are integrable discrete series of
\(H_\infty\) and \(G_\infty\), respectively, with
\[
 \sigma\lhook\joinrel\longrightarrow\pi|_{H_\infty},
 \qquad h>4n^2+2
\]
where \(h\) is a lower bound for the real opposite-sign gaps of \(\pi\)
in Proposition~\ref{prop:bessel-sign-gap-decay}. Write
\([H]=H(F)\backslash H(\mathbb A_F)\).

One finite depth \(d_*\) suffices for every nonzero smooth row vector
\(u\in\sigma^\infty\) and every nonzero cuspidal multiplicity map
\(t\) of type \(\sigma\) and level \(J_H\). Some cusp form \(P\)
of real type \(\pi\), right invariant by \(K^vK_v(d_*)\), satisfies
\[
                 \int_{[H]}P(h)\overline{t(u)(h)}\,dh\ne0 .
\]
For the bilinear period take \(t(u)\) to be the conjugate of the row
form. At each real place where the subgroup fixes a positive line,
He's signed interlacing criterion determines the displayed branching
condition. \cite[Theorem~0.1]{he-discrete-branch}
\end{auxlemma}
\begin{proof}
\proofstep{Choose a coefficient with a nonzero convolution on the row.}
First fix a compact-finite unit vector \(u\in\sigma\) and an isometric inclusion
\(j:\sigma\to\pi|_{H_\infty}\). Use inner products linear in the first
entry. Schur orthogonality makes
\(d_\sigma\overline{\langle\sigma(h)u,u\rangle}\) a self-adjoint
convolution idempotent, where \(d_\sigma\) is the formal degree.
It is in \(L^1(H_\infty)\) by integrability. Its image in \(\pi\)
contains \(ju\). Choose a compact-finite \(b\) with
\(\langle b,ju\rangle\ne0\) and set
\[
 k(g)=\langle\pi(g)b,b\rangle,\qquad
 A=\int_{H_\infty}\overline{k(h)}\,\sigma(h)\,dh .
\]
The gap estimate makes the second integral absolutely convergent.
Pairing the idempotent applied to \(b\) with \(b\) gives
\(d_\sigma\langle Au,u\rangle>0\).

Realize \(t\) on the Hilbert space of cusp forms with its Petersson
inner product. As \(Au\ne0\), injectivity gives \(t(Au)\ne0\).
For \(f=t(u)\), its value in \(L^2\) is the convolution of \(f\)
by \(\overline{k}|_{H_\infty}\). Thus for some \(x\in H(\mathbb A_F)\),
\[
                 \int_{H_\infty} k(h)\overline{f(xh)}\,dh\ne0 .
\]
The integral is pointwise defined since the smooth compact-finite cusp
form \(f\) is bounded.

\proofstep{Keep the stabilizer volume fixed while shrinking the level.}
Put \(J=J_{H,v}\cap K_v(0)\), where \(K_v(0)\) is the maximal compact
of the frame. The subgroups \(J K_v(d)\) are open because \(K_v(d)\)
is normal in \(K_v(0)\); their intersection is \(J\). Take the
Poincar\'e series
\[
 P_{x,d}(g)=\sum_{\gamma\in G(F)}
      k(x_\infty^{-1}\gamma_\infty g_\infty)
      {\bf1}_{K^v J K_v(d)}(x_f^{-1}\gamma_f g_f).
\]
These series belong to the cuspidal real type \(\pi\). Indeed right
differentiation and compact translation act on the vector in the first slot of
the coefficient. Unfolding a proper constant term gives zero by the
vanishing of the unipotent integrals of the discrete-series
coefficients. This is the integrable discrete-series construction
of cuspidal Poincar\'e series. \cite{muic-fourier-poincare}

Here one absolute bound works for every \(d\) and \(g\). For a fixed
\(x\) all distinct surviving real arguments have quotients in the
fixed discrete subgroup obtained from \(G(F)\) by the finite cutoff
at \(d=0\) and conjugation by \(x\). In a product of boxes of matrix
norm \(R_w\ge1\), a neighborhood of the identity separating these
quotients gives a packing bound
\[
           \#\{\text{arguments in the boxes}\}
                      \ll_x\prod_{w\mid\infty}(1+R_w)^{4n^2}.
\]
In orthonormal signature frames \(\|b^{-1}\|=\|b\|\) for the
Hilbert--Schmidt norm. Thus \(\|ab^{-1}-1\|\le R_w\|a-b\|\) there, and a complex
\(n\)-by-\(n\) matrix has \(2n^2\) real coordinates. Apply
Proposition~\ref{prop:bessel-sign-gap-decay} to a finite tensor
expansion of the coefficient with exponent \(4n^2+1\). Dyadic
summation gives the bound, also for each fixed derivative.

\proofstep{Unfold the limiting period.}
For \(h\in H(\mathbb A_F)\), a rational \(\gamma\) surviving at all
depths satisfies \(x_v^{-1}\gamma_v h_v\in J\subset H(F_v)\).
The equation fixing the distinguished line vector therefore holds over
\(F\) for \(\gamma\). Thus precisely the surviving terms of \(H(F)\)
contribute in the limit. The preceding absolute bound and finite volume
of \([H]\) justify dominated convergence. Unfolding the surviving sum
gives
\[
 \lim_{d\to\infty}\int_{[H]}P_{x,d}(h)\overline{f(h)}\,dh
   =\operatorname{vol}\!\left((K^v\cap H(\mathbb A_f^v))J\right)
                \int_{H_\infty} k(h)\overline{f(xh)}\,dh .
\]
The finite volume factor is positive since its group is open in
\(H(\mathbb A_f)\), and it lies in \(J_H\). The right side is
nonzero, so for this \(t\) some finite depth works.

Finally the row multiplicity space at the fixed type and level is
finite-dimensional. Its common kernels under all such period tests
at levels \(K^vK_v(d)\) form a decreasing sequence. Their intersection
is zero by the preceding argument, so one of them is zero. This gives
\(d_*\) for the chosen compact-finite \(u\).

\proofstep{Keep that depth for every smooth row vector.}
For each fixed nonzero \(t\), choose its period test \(P\) at
\(d_*\) by the preceding construction. The functional
\[
 \ell_{P,t}(v)=\int_{[H]}P(h)\overline{t(v)(h)}\,dh,
 \qquad |\ell_{P,t}(v)|\le\|P|_{[H]}\|_2\,\|t\|\,\|v\|
\]
is continuous on the Hilbert space of \(\sigma\), since this \(P\)
is bounded and \([H]\) has finite volume. It is nonzero on \(u\).
For any nonzero smooth \(v\), irreducibility makes the linear span
of \(\sigma(h_0)v\), \(h_0\in H_\infty\), dense. Choose \(h_0\)
on which the functional is nonzero and move the right translation
to \(P\) by invariance of the quotient measure:
\[
 \ell_{P,t}(\sigma(h_0)v)=\ell_{R(h_0^{-1})P,t}(v).
\]

This translated test keeps its finite level. To take a compact-finite
test there, convolve it on the right with probability approximate
identities of finite Peter--Weyl support on the real maximal compact
group. The translated cusp form is bounded and continuous, so these
convolutions converge pointwise, with a common bound. They remain in
real type \(\pi\), and their periods converge by dominated convergence
since \(t(v)\in L^1([H])\). Some compact-finite test has nonzero period.
\end{proof}

The lemma first fixes all spaces and weights, then finds a depth. In
Definition~\ref{hyp:fitted-pointwise-nonvanishing} the depth is
prescribed before the vector is tested and the spaces vary with the
scale. Dominated convergence at a fixed datum gives no comparison
of \(d_*\) with that prescribed depth. A bottom form obtained here
also has to be in the span of the descending columns at the required
level before the next Fourier--Jacobi or Bessel step can use it. Lemmas~\ref{lem:fitted-prescribed-tail} and~\ref{lem:fitted-core-span} supply these two comparisons.

For the prescribed level put
\[
 f=[F:\mathbb Q],\qquad r=\dim_EY_0,\qquad n=r+1,
 \qquad L=\log b.
\]
Write \(g_0,g_1,p_d,p'_d\) for the two level primes
and two auxiliary primes of the datum. Over the least split auxiliary
rational prime \(q_0\) choose the place \(v_0\) of the grid. Let
\(S_0\) comprise \(v_0\) and the places over \(g_0g_1p_dp'_d\).
For each \(v\in S_0\) use the specified prime \(\mathfrak w_v\) of
\(E\) on one split factor. The companion level uses the ideal
\begin{equation}\label{eq:fitted-auxiliary-depth}
 \begin{split}
 D&=C_F f^2 n^6\log\!\bigl(b n g_0g_1p_dp'_d|\operatorname{disc}E|\bigr),\\
 d_{v_0}&=1+\left\lfloor D/\log q_0\right\rfloor,\qquad
 d_v=2\quad(v\in S_0\setminus\{v_0\}),\qquad
 \mathfrak m=\prod_{v\in S_0}\mathfrak w_v^{d_v} .
 \end{split}
\end{equation}
The floor determines a strict lower bound even when \(D/\log q_0\)
is an integer. Indeed \(\mathrm N\mathfrak w_{v_0}\ge q_0\), so
\[
 \log\mathrm N\mathfrak m\ \ge\ d_{v_0}\log q_0
           \ >\ D\ \ge\ C_F f^2 n^6 L .
\]

\begin{auxlemma}[Numerical allowance at the prescribed level]\label{lem:fitted-depth-allowance}
Use \eqref{eq:fitted-auxiliary-depth} with \(C_F\ge6510\). Let
\(1\le m,f\le n\), \(L\ge1\), \(C\ge0\), \(n\ge C+1\), and
\(0\le Q\le C f n^4L\). Here \(m\) is the chain length. Put
\(A=r^2 f g_0g_1|\operatorname{disc}E|^{1/(2f)}\).
Assume the grid size bounds \(\log(g_0g_1)\le20L\) and
\(\log|\operatorname{disc}E|^{1/(2f)}\le L\), with \(r\ge1\).
For \(s=\tfrac12(2e^Q A)^{-m}\) and
\(R=s^f\sqrt{\mathrm N\mathfrak m}\),
\[
 0<s\le\tfrac12,\qquad e^Q A(2s)^{1/m}=\tfrac12,\qquad
 R>2e^{Q+(6n+9)f} .
\]
\end{auxlemma}
\begin{proof}
We have \(A\ge1\) and \(\log A\le n^3+21L\) since \(r^2 f\le n^3\).
The first two assertions follow from the definition of \(s\). For the
last one compute
\[
 \log 2+Q+(6n+9)f+f\log(1/s)
 = (1+f+fm)\log 2+(1+fm)Q+(6n+9)f+fm\log A .
\]
Use \(m\le n\) and \(\log 2<1\). The terms with \(C\) are bounded by
\(2C f^2 n^5 L\). The remaining terms are at most
\((3+15+1+21)f^2 n^5 L\). As \(n\ge C+1\), the sum is at most
\(3255 f^2 n^6 L < \tfrac12\log\mathrm N\mathfrak m\).
Exponentiation gives the stated strict inequality.
\end{proof}

\begin{auxlemma}[Lattice lengths at a large value of the row]\label{lem:fitted-large-row-lattice}
Use the base row \(Y_0\), its self-dual Hermitian lattice
\(\Lambda_Y\), the principal row level of modulus \(g=g_0g_1\), and
the parameters of Lemma~\ref{lem:fitted-depth-allowance}.
Let \(f_0\ne0\) be a smooth cuspidal form invariant under that level.
In the real signature frames assume the directional bound
\[
 \|R(X)f_0\|_\infty\le e^Q\|f_0\|_\infty\|X\|_{\rm HS}
 \quad\text{for }X\text{ in any one real row Lie algebra}.
\]
Here \(R(X)\) denotes right differentiation. For
\(|f_0(x)|>\|f_0\|_\infty/2\) put
\(L_x=Y_0(E)\cap x_f\widehat\Lambda_Y\) with positive real norms
\(\|y\|_{x,\tau}=\|x_\tau^{-1}y_\tau\|\). For every nonzero
integral ideal \(\mathfrak a\) one has
\[
 \prod_{\tau\mid\infty}\|y\|_{x,\tau}
       \ge s^f\sqrt{\mathrm N\mathfrak a}
       \qquad(0\ne y\in\mathfrak a L_x).
\]
\end{auxlemma}
\begin{proof}
The transported lattice is self-dual since each \(x_v\) is unitary.
Its covolume sizes and principal level are those used in the covering
argument of Proposition~\ref{prop:short-orbit-cover}. This geometric
part of that proof uses self-duality, the integral trace-one element
of the finite-unramified extension \(E/F\), and the negative index
as an upper bound on the rank of an isotropic flag. The negative
index of \(Y_0\) is \(m/2\); the weaker upper bound \(m\) suffices.

Write \(t\) for the smallest filtration size of \(L_x\). If
\(t<s\le1/2\), the covering gives real representatives on a rational
parabolic unipotent orbit through \(x\) of the form \(\exp U\exp V\).
Use Cauchy--Schwarz to sum the real-place norms in the proposition:
\[
 \sum_\tau(\|U_\tau\|_{\rm HS}+\|V_\tau\|_{\rm HS})
 \le \tfrac14 r^2 f g\,|\operatorname{disc}E|^{1/(2f)}(2t)^{1/m}
 =\tfrac14 A(2t)^{1/m} .
\]
Integrating the directional bound along these segments controls the
variation from \(f_0(x)\). Average on the compact unipotent quotient
with probability measure. Its cusp average is zero, so
\[
 |f_0(x)|\le \tfrac14e^Q A(2t)^{1/m}\|f_0\|_\infty
              <\tfrac12\|f_0\|_\infty .
\]
This contradicts the choice of \(x\), giving \(t\ge s\).

For \(y\in\mathfrak aL_x\setminus\{0\}\) write the saturated line
\(L_x\cap Ey=\mathfrak b y\) with \(\mathfrak b\) a fractional ideal.
Projectivity gives \(\mathfrak aL_x\cap Ey=\mathfrak a\mathfrak b y\),
so \(1\in\mathfrak a\mathfrak b\) and
\(\mathrm N\mathfrak b\le(\mathrm N\mathfrak a)^{-1}\).
The rank-one covolume, with one embedding per conjugate pair, gives
\[
 \prod_\tau\|y\|_{x,\tau}
 =\operatorname{size}(\mathfrak b y)^f(\mathrm N\mathfrak b)^{-1/2}
 \ge t^f\sqrt{\mathrm N\mathfrak a} .
\]
\end{proof}

\begin{auxlemma}[The tail of the period at the auxiliary level]\label{lem:fitted-prescribed-tail}
Use Lemma~\ref{lem:fitted-large-row-lattice}, the inclusion
\(\iota:H=U(Y_0)\hookrightarrow G=U(C_0)\), and the norm-one vector
\(z_0\) in the added line. Let \(K_C\) be the principal subgroup on
\(\Lambda_Y\perp\mathcal O_E z_0\) for the ideal \(\mathfrak m\) in
\eqref{eq:fitted-auxiliary-depth}. Let \(k\) be a real matrix
coefficient of the required ambient representation for which
Poincar\'e summation gives a cusp form. For the row \(f_0\), at the
point \(x\) of that lemma, suppose a common scalar \(c\ne0\) satisfies
\[
 \begin{aligned}
 \int_{H_\infty}k(\iota h)\overline{f_0(xh)}\,dh
  &=c\overline{f_0(x)},\\
 \int_{H_\infty}|k(a\iota h)|\,dh
  &\le |c|e^Q\prod_\tau\|a_\tau z_0\|^{-(2n+5)}
      \qquad(a\in G_\infty).
 \end{aligned}
\]
The period of the Poincar\'e series with seed \(k\,{\bf1}_{K_C}\),
left translated by \(\iota x\), against \(f_0\) is nonzero.
\end{auxlemma}
\begin{proof}
Unfold by right \(H(F)\)-cosets in \(G(F)\). For each supported coset the
finite integral has volume \(V=\operatorname{vol}(K_C\cap H(\mathbb A_f))\).
This compact group fixes \(f_0\): congruence modulo \(\mathfrak m\) and
self-duality also imply congruence modulo \(\overline{\mathfrak m}\)
by taking inverse adjoints. Together the ideals impose the full row
level at \(g_0g_1\). The stabilizer coset thus gives \(V c\overline{f_0(x)}\).

A coset \(\gamma H(F)\) is determined by \(\gamma z_0\). The finite
support sends \(y=\gamma z_0-z_0\) into
\(\mathfrak m(L_x\perp\mathcal O_Ez_0)\). Distinct cosets give distinct
vectors. For a nonzero vector of this lattice its row component or its
line component is nonzero, so Lemma~\ref{lem:fitted-large-row-lattice}
or the ideal norm on the line gives
\[
    \prod_\tau\|y\|_{x,\tau}\ge R=s^f\sqrt{\mathrm N\mathfrak m}>1 .
\]

The transported unit vector \(\gamma z_0\) has positive norm at least
one in each signature frame. Hence
\(1+\|y\|_{x,\tau}\le3\|\gamma z_0\|_{x,\tau}\).
The transversal integral in the hypothesis now bounds the tail by
\(V |c|e^Q\|f_0\|_\infty\) times
\[
     3^{(2n+5)f}\sum_{y\ne0}\prod_\tau(1+\|y\|_{x,\tau})^{-(2n+5)} .
\]

Here the sum is over the whole indicated lattice. For the box
\(2^{j_\tau}\le1+\|y\|_{x,\tau}<2^{j_\tau+1}\), put
\(D=2^{f+\sum_\tau j_\tau}\). Scale each factor by
\(2^{-j_\tau-1}\). In the maximum norm the box lies in the unit ball,
and the product lower bound on differences gives separation at least
\(D^{-1/f}\). Packing in real dimension \(2nf\) gives at most
\((1+2D^{1/f})^{2nf}\le3^{2nf}D^{2n}\) points.

Take one power out of the product using \(R^{-1}\). Summing the
remaining box bounds gives the convergent majorant
\[
 \frac{3^{(4n+5)f}2^{2nf}}{R}
     \prod_\tau\sum_{j\ge0}2^{-4j}
 =\frac{(3^{4n+5}2^{2n}\,16/15)^f}{R}
 \le\frac{e^{(6n+9)f}}{R} .
\]
This justifies the absolute unfolding. Lemma~\ref{lem:fitted-depth-allowance}
makes the absolute tail smaller than \(V|c|\|f_0\|_\infty/2\), which is
strictly below the stabilizer term in absolute value.
\end{proof}

\begin{auxlemma}[A bounded translate meeting the prescribed row]\label{lem:weighted-fitted-translate}
Let \(z_0=e_{i_0}\) be a positive coordinate vector in the signature
frame of \(G=U(p,q)\), let \(H=G_{z_0}\), and put \(n=p+q\ge4\).
Write \(\lambda,\mu\) for the decreasing signed Harish--Chandra
parameters of a discrete series \(\pi\) of \(G\) and a
multiplicity-one discrete-series summand \(j:\sigma\hookrightarrow
\pi|_H\), respectively, with \(j\) isometric. Assume
\[
 \min_{i<n}(\lambda_i-\lambda_{i+1})\ge2,\qquad
 |\lambda_i|,|\mu_j|\le B=b^2,\qquad b\ge n.
\]
Assume that entries of opposite signatures in each parameter differ
by at least \(1000n^3\). For every unit vector \(v_0\) in the lowest
compact type of \(\pi\), some \(a\in G\) satisfies, for \(T=8nB\),
\[
 \|a\|_{\rm HS}\le (n+4T^{8n})^{1/2}\le T^{10n},\qquad
 \|j^*\pi(a)v_0\|\ge T^{-C n^3}
\]
with an absolute \(C\).
\end{auxlemma}
\begin{proof}
Take a unit vector \(u_0\) in the lowest compact type of \(\sigma\)
and put \(P=j^*,\ \xi=ju_0\). For the given inner products, linear in
their first entry, put \(F(g)=\langle\pi(g)v_0,\xi\rangle\).
We will localize its square integral transversally to \(H\) by proving
\begin{equation}\label{eq:weighted-fitted-moment}
 \int_G\|g^{-1}z_0\|^2|F(g)|^2\,dg
       \le T^{8n}/d_\pi,\qquad
 \int_G|F(g)|^2\,dg=1/d_\pi .
\end{equation}

\proofstep{Translation by the dual standard representation.}
Write \(x_l(y)=y\otimes e_l^*\) and
\(\Pi=\pi^\infty\otimes(\mathbb C^n)^\vee\).
Filter the dual standard module by the flag for the upper triangular
Borel in the ordered signature frame. The graded characters have
weights \(-e_i\). The tensor identity for cohomological induction
and good-range concentration for the compact Cartan give the graded
modules
\[
       \pi_i=\pi_{\lambda-e_i},\qquad 1\le i\le n .
\]
Indeed every shifted parameter is still strictly decreasing, and
each Borel induction occurs only in degree
\(\dim(\mathfrak n\cap\mathfrak k)\). The resulting short exact
sequences split into distinct infinitesimal-character summands.
Smooth globalization gives the same decomposition for \(\Pi\).
\cite[Proposition~3.6 and (14)]{prudhon-translation}
\cite[Theorems~5.35, 5.99, 8.2 and Chapter~XI]{knapp-vogan-cohomological}
\cite[Theorems~1.1 and 10.6]{casselman-wallach}

Let \(p_i\) denote these projections, and let \(\epsilon_l\) be the sign
of the \(l\)-th coordinate. The total Casimir is symmetric for the
invariant Hermitian form \([x,y]=\sum_l\epsilon_l\langle x_l,y_l\rangle\).
Its real eigenvalues are distinct. Thus the summands are orthogonal
and the form on each is nondegenerate. Uniqueness of the invariant
Hermitian form on the irreducible module gives quotient maps \(R_i\)
to unitary models, normalized by
\[
         [p_i x,y]=s_i\langle R_i x,R_i y\rangle,\qquad s_i=\pm1 .
\]

\proofstep{The coefficient on the fixed coordinate.}
Set \(\mathcal E=(d\pi(E_{ab}))_{ab}\), with \(E_{ab}\) the standard
matrix units. Expanding the quadratic Casimir on the dual tensor gives
\[
 C_\Pi=C_\pi I+nI-2\mathcal E,\qquad
 \mathcal E p_i=(\lambda_i+(n-1)/2)p_i .
\]
Let \(Q_n(u)\) be the Capelli column determinant with successive
shifts \(u,u-1,\ldots,u-n+1\). In the adjugate \(A(u)\), entry
\((l,a)\) deletes row \(a\) and column \(l\), includes the cofactor
sign, and uses consecutive shifts starting at \(u\). The alternating
column expansion with the last column distinguished gives
\[
 A(u)(\mathcal E+(u-n+1)I)=Q_n(u)I,\qquad
 Q_n(u)|_\pi=\prod_k(u+\lambda_k-(n-1)/2).
\]
For the first identity specialize the quantum determinant to
\(T(u)=I+\mathcal E/u\). Its alternating expansion vanishes for
repeated columns. For the second identity evaluate the central
element on a highest-weight module of the same infinitesimal
character, where its diagonal product acts on the highest-weight
line. \cite[Proposition~2.7(4) and Remarks~2.8(1), 2.11]{mno-yangians}

Divide the first identity on each summand by its scalar linear
factor, then evaluate at \(u_i=(n-1)/2-\lambda_i\). Deleting the
fixed coordinate gives a rank-\(n-1\) Capelli determinant. Since
\(P\) intertwines the derived \(H\)-action on smooth vectors, we obtain
\begin{align*}
 p_i&=A(u_i)\Big/\prod_{k\ne i}(\lambda_k-\lambda_i),\\
 P\big((p_i x_{i_0}(y))_{i_0}\big)&=\gamma_i Py,\qquad
 \gamma_i=\frac{\prod_j(\mu_j-\lambda_i+1/2)}
                      {\prod_{k\ne i}(\lambda_k-\lambda_i)} .
\end{align*}
Each denominator has absolute value at least one. For words of
length at most \(n\), \eqref{eq:count-step} bounds each derived
factor on the preceding filtered space by \(T^2\). Expanding the
\((l,l)\) minor thus gives
\[
 |\gamma_i|\le T^n,\qquad
 \|R_i x_l(v_0)\|^2
  =|[p_i x_l(v_0),x_l(v_0)]|
  \le(n-1)!(B+n+T^2)^{n-1}\le T^{4n}.
\]

\proofstep{Schur orthogonality with the nonsmooth ambient vector.}
The matrix unit
\[ e(h)=d_\sigma\overline{\langle\sigma(h)u_0,u_0\rangle} \]
is integrable by \eqref{eq:count-polar-decay} and the ball estimate
below. Its self-adjoint convolution idempotent acts on \(\pi|_H\)
as projection onto \(\mathbb C\xi\), by multiplicity one: a second
nonzero image on the orthogonal complement would generate a second
copy of \(\sigma\) by the matrix-unit identities.

Denote the
orthogonal projection it defines on \(\pi_i|_H\) by \(Q_i'\).
Choose smooth \(\eta\) with \(\langle\eta,\xi\rangle=1\) and set
\(b_i=Q_i'R_i x_{i_0}(\eta)\). Averaging the invariant form against
\(e\), using that \(H\) fixes the positive coordinate, gives
\[
 \langle(p_i y)_{i_0},\xi\rangle
    =s_i\langle R_i y,b_i\rangle,\qquad
 s_i\|b_i\|^2=\gamma_i .
\]
The second equality uses \(y=x_{i_0}(\eta)\). In particular no
ambient derivative of \(\xi\) occurs. For each \(l\),
\[
  (g^{-1})_{li_0} F(g)
   =\sum_i s_i\langle\pi_i(g)R_i x_l(v_0),b_i\rangle .
\]
The formal-degree formula for the same group and Haar measure gives
\[
 d_\pi/d_{\pi_i}
  =\prod_{k\ne i}\left|\frac{\lambda_i-\lambda_k}
                                  {\lambda_i-1-\lambda_k}\right|
  \le2^{n-1}.
\]

The compact central factor has the same measure in both degrees.
\cite[Theorem~8.1 and (3.10), (3.12)]{atiyah-schmid}
Schur orthogonality, Cauchy--Schwarz in \(i\), and summation in \(l\)
give the explicit moment estimate
\[
 d_\pi\int_G\|g^{-1}z_0\|^2|F(g)|^2\,dg
          \le n^3T^{5n}2^{n-1}\le T^{8n}.
\]
Ordinary Schur orthogonality gives the unweighted equality.

\proofstep{A compact set of translates.}
Put \(S=\{g:\|g^{-1}z_0\|^2\le2T^{8n}\}\). By
\eqref{eq:weighted-fitted-moment} this left \(H\)-invariant set
carries square integral at least \(1/(2d_\pi)\). For
\(x=g^{-1}z_0\) write \(x=x_++x_-\) in the positive and negative
subspaces, so \(\|x_+\|^2-\|x_-\|^2=1\). A compact rotation followed
by the real hyperbolic block with diagonal \(\|x_+\|\) and
off-diagonal \(\|x_-\|\) carries \(z_0\) to \(x\).

Its inverse represents \(Hg\) and has Hilbert--Schmidt norm at most
\[
       R=(n+4T^{8n})^{1/2}.
\]
Write \(B_D(t)=\{h\in D:\|h-1\|_{\rm HS}<t\}\),
\(D(L)=\{h\in D:\|h\|_{\rm HS}\le L\}\) and
\(V_H=\operatorname{vol} B_H(1/4)\).
Over every such coset the translates by \(B_H(1/4)\) of its
representative lie in \(G(2R)\). Quotient integration and Schur
orthogonality on \(H\) therefore give
\begin{align*}
 \operatorname{vol}(H\backslash S)&\le\operatorname{vol}G(2R)/V_H,\\
 \int_H|F(hg)|^2\,dh&=\|P\pi(g)v_0\|^2/d_\sigma .
\end{align*}
A coset with at least half the average has a representative \(a\)
of norm at most \(R\) satisfying
\[
 \|P\pi(a)v_0\|^2
      \ge d_\sigma V_H/(4d_\pi\operatorname{vol}G(2R)).
\]

For either matrix group \(D\) of size \(r\), the identity
\(h^{-1}=Jh^*J\) and Euclidean packing of a maximal
\(t/L\)-separated subset of \(D(L)\) give
\[
 \operatorname{vol}D(L)\le(1+2L^2/t)^{2r^2}\operatorname{vol}B_D(t).
\]
Take \(t=(4T^2)^{-1}\). On \(B_G(t)\) the matrix logarithm has norm
at most \(2t\). The first derivative bound \eqref{eq:count-step}
shows that the lowest-type diagonal coefficient of \(v_0\) has
modulus at least \(1/2\) there. Schur and packing give
\[
 d_\pi\operatorname{vol} B_G(t)\le4,\qquad
 d_\pi\operatorname{vol}G(2R)\le T^{C_1 n^3}.
\]

It remains to bound \(d_\sigma V_H\) from below in the same Haar
measure. Applied to the row itself, \eqref{eq:count-polar-decay}
bounds its lowest-type diagonal coefficient by
\(e^{C_2 n^2\log n}\|h\|_{\rm HS}^{-40(n-1)^2}\):
for Cartan length at least one use the opposite-sign gap, and below
one use unitarity and \(\|h\|_{\rm HS}\le e\sqrt n\).
The same assertion holds for a compact row by unitarity.
Summing the packing estimate on dyadic balls now gives
\[
  1/d_\sigma=\int_H|\langle\sigma(h)u_0,u_0\rangle|^2\,dh
       \le V_H\,e^{C_3 n^2\log n}.
\]
Together the last three inequalities give the claimed bound for
the projection of the same \(a\).
\end{proof}

\begin{auxlemma}[Normalization by a projection on the row type]\label{lem:fitted-real-projection}
At one real place let \(G=U(p,q)\), \(H=U(p-1,q)\), \(n=p+q\ge4\).
Let \(j:\sigma\hookrightarrow\pi|_H\) be a multiplicity-one Hilbert
discrete-series summand. Suppose the two parameters have coordinates
bounded by \(b^2\), with \(b\ge n\), and opposite-sign gaps at least
\(1000n^3\). Let \(v_0\) be a unit vector in the lowest compact type
of \(\pi\). Put \(T=8nb^2\). Suppose there is a translate \(v=\pi(a)v_0\)
for which the same matrix \(a\) satisfies
\begin{equation}\label{eq:fitted-real-projection}
       \|a\|_{\rm HS}\le T^{10n},\qquad
       \|j^*v\|\ge T^{-C_0 n^4}
\end{equation}
with a fixed \(C_0\ge0\).

Write \(w=j^*v/\|j^*v\|\) and
\(c=\|j^*v\|^2/d_\sigma\) for the row formal degree \(d_\sigma\).
For \(k(g)=\langle\pi(g)v,v\rangle\) and inner products linear in
the first entry, the row convolution by \(\bar k\) sends \(w\) to \(cw\).
With \(Q_1=C_1(1+C_0)n^4\log T\) for an absolute \(C_1\),
\[
 \begin{aligned}\int_H|k(gh)|\,dh&\le ce^{Q_1}\|gz_0\|^{-(2n+5)},\\
 \int_H\left|\left.\tfrac{d}{dt} k(\exp(tX)h)\right|_{t=0}\right|\,dh
    &\le ce^{Q_1}\quad (\|X\|_{\rm HS}\le1,\ X\in\mathfrak h).\end{aligned}
\]
Here \(z_0\) is the fixed positive vector of norm one.
\end{auxlemma}
\begin{proof}
For compact-finite row vectors, integrable discrete-series matrix
units act on \(j\sigma\) by Schur orthogonality. They vanish on its
orthogonal complement: the translates of the image of a nonzero
matrix-unit projection generate copies of \(\sigma\), by the same
matrix-unit multiplication rules. Multiplicity one excludes such a
copy there. Pairing these identities against \(v\) shows that the
Fourier coefficient of \(\bar k|_H\) in \(\sigma\) is the operator
\[
    \int_H\bar k(h)\sigma(h)\,dh
    =d_\sigma^{-1}\langle\,\cdot\,,j^*v\rangle j^*v .
\]
The coefficient is integrable by the decay below, so the operator
identity extends by continuity. As \(v\) is smooth and \(j^*\) is a
bounded intertwiner, \(w\) is smooth for \(H\).

\proofstep{Numerators in the Haar measure of the row.}
Write \(D=40n^2\). For a Cartan factor \(\exp X\), put
\(t=\|X\|_{\rm op}\). When \(t\ge1\), the polar estimate
\eqref{eq:count-polar-decay} gives \(e^{-Dt}\), since
\(h/(8\sqrt{2q})\ge D\). For \(t<1\) use unitarity, also covering
the compact case. The signature-frame norm is at most \(\sqrt n e^t\).
Conjugation by \(a\) in \eqref{eq:fitted-real-projection} therefore gives
\[
           |k(g)|\le M\|g\|_{\rm HS}^{-D},
           \qquad \log M\ll n^3\log T .
\]

Put \(m_1=n-1\), \(u=4m_1^2+1\), and \(V_H=\mathrm{vol}\{h:\|h-1\|_{\rm HS}<1/4\}\).
In the box \(\|h\|\le R\), left translates of this ball centred on
a maximal \((4R)^{-1}\)-separated set give a cover, as
\(\|h^{-1}\|=\|h\|\). Euclidean packing gives at most
\((1+8R^2)^{2m_1^2}\) translates. Summing on dyadic boxes yields
\[
                  \int_H\|h\|_{\rm HS}^{-u}\,dh\le e^{O(n^2)}V_H .
\]

Write \(g=b_g h_0\) with \(h_0\in H\) and
\(\|b_g\|_{\rm HS}\le\sqrt{n+2}\,\ell\), \(\ell=\|gz_0\|\ge1\).
To build \(b_g\), rotate compactly and apply a boost on the plane of
the positive and negative parts of \(gz_0\); its diagonal and
off-diagonal entries are the lengths of those parts. Thus
\(\|b_gh\|\ge\ell\) and \(\|h\|\le\sqrt{n+2}\,\ell\|b_gh\|\). It follows that
\[
 \int_H |k(gh)|\,dh\le M(n+2)^{u/2}\ell^{-D+2u}\int_H\|h\|^{-u}\,dh,
       \qquad D-2u\ge 2n+5 .
\]

For \(F_h(t)=k(\exp(tX)h)\) and \(|t|\le1\) the decay bound gives
\(|F_h(t)|\le M e^D\|h\|^{-D}\). By the two Casimir word bounds
\eqref{eq:count-step}, differentiating the second slot twice gives
\(\|F_h''\|_\infty\le b_0 b_1\|a\|_{\rm HS}^4\), where \(b_0,b_1\) are those
bounds in dimension \(n\). Taylor's estimate at \(t=\|h\|^{-D/2}\) yields
\[
 |F'_h(0)|\le (2Me^D+b_0b_1\|a\|_{\rm HS}^4)\|h\|^{-D/2} .
\]
Since \(D/2\ge u\), the integrated first derivative costs at most
\(V_H\exp(O(n^3\log T))\), as does the transversal bound without its
power of \(\ell\).

\proofstep{The same denominator for both estimates.}
For a lowest-type row unit vector the first Casimir bound in dimension
\(n-1\) is \(B=\sqrt{n-1}\,b^2+\tfrac32(n-1)^{3/2}\).
Its diagonal coefficient has absolute value at least \(1/2\) on
\(\|h-1\|_{\rm HS}\le\delta=1/(4B)\). Indeed the matrix logarithm
has norm at most \(2\delta\) there, so integration along its exponential
changes the vector by at most \(2B\delta\). Schur orthogonality gives
\(d_\sigma\operatorname{vol}(B_H(\delta))\le4\), where \(B_H(\delta)\)
denotes this ball. By the same packing bound, at radius
\(\|h\|_{\rm HS}\le\sqrt{n-1}+1\),
\[
 d_\sigma V_H\le4(1+2(\sqrt{n-1}+1)^2/\delta)^{2(n-1)^2},
 \qquad V_H/c\le T^{O((1+C_0)n^4)} .
\]
Dividing the preceding two bounds by \(c\) proves the lemma.
\end{proof}

Tensor the coefficients over the \(f\) real places. The reproducing
scalar is \(\prod_\tau c_\tau\). Lemma~\ref{lem:weighted-fitted-translate}
gives \eqref{eq:fitted-real-projection} at every place once the
ordered bottom parameters have adjacent gaps at least two. The
large-scale bottom lists satisfy this condition: their consecutive
middle entries satisfy \(v_i-v_{i+1}\ge2\), their lower joining gap
is \(v_k+b_1\ge2\), and the outer blocks have the larger stipulated
separations. These are conditions on the bottom pair. Later spherical
transfers insert consecutive strings of a different length.

Both normalized integral estimates now have the same bound
\[
          Q\ll f n^4\log b .
\]
Differentiating the reproducing identity moves the derivative onto
the kernel. Lemma~\ref{lem:fitted-large-row-lattice} applies to the
resulting derivative estimate, and Lemma~\ref{lem:fitted-prescribed-tail}
bounds the terms off the stabilizer at the prescribed level. To
use this bottom period on the tower, one must next realize its test
as a linear combination of descents from permitted columns.

\begin{auxlemma}[A common split test for a whole space of rows]\label{lem:fitted-split-row-projector}
At a marked split place of a core arrow
\(C_i\to C_{i+1}\), let \(O\) be the ring of integers of the chosen
half, \(I\subsetneq O\) its prescribed depth ideal, and
\(X=\operatorname{Hom}_O(\Lambda_{i+1},\Lambda_i)\) in the two
lattice frames. Put \(d=\operatorname{rank}_O\Lambda_i\) and
\(J=\ker(\mathrm{GL}_d(O)\to\mathrm{GL}_d(O/I))\).
Let \(x_0\in X\) be surjective modulo \(I\). If the splitting
character is trivial on \(1+I\), the local Weil representation
on the chosen half satisfies
\[
 f=\mathbf1_{x_0+IX},\qquad
 \langle\omega(g)f,f\rangle=\|f\|_2^2\,\mathbf1_J(g)
       \quad (g\in\mathrm{GL}_d(F_v)).
\]
At all marked split places these functions can be chosen using a
single integral homomorphism \(x_0\).
\end{auxlemma}
\begin{proof}
Choose a projection in each pair of residue frames. Reduction of the
homomorphism module commutes with base change since the source lattice
is projective. The Chinese remainder theorem therefore lifts the
projections simultaneously. If \(X^\vee\) denotes the character-dual
lattice, finite character orthogonality realizes the result in the
allowed depth space:
\[
 f(x)=\mathbf1_X(x)\,
       \frac{1}{[I^{-1}X^\vee:X^\vee]}
       \sum_{\lambda\in I^{-1}X^\vee/X^\vee}
                         \psi_v(\langle x-x_0,\lambda\rangle).
\]
The split action sends the support to \(g(x_0+IX)\). Suppose the
two supports meet, so \(a=gb\) for \(a,b\in x_0+IX\).
Choose an integral \(y\) with \(x_0 y\equiv 1\pmod I\). Then
\[
     ay,by\in 1+I M_d(O),\qquad g=(ay)(by)^{-1}\in J .
\]
Conversely, \(J\) preserves the coset. Its determinants have absolute
value one and the splitting scalar is one. The other half acts by
inverse adjoint, including at a split dyadic place. This proves the
formula with the original local scalar product.
\end{proof}

\begin{auxlemma}[The spherical factor on a core chain]\label{lem:fitted-core-euler}
Put \(n=r+1\), \(h=r+5\) and \(N_i=n+ih\).
Denote by \(S\) the archimedean and exceptional places. Start with a
spherical standard label \(B=(B_w)_{w\nmid S}\) of degree \(n\)
over the CM field \(E\), with
\(\max_{\alpha\in B_w}|\alpha|\le q_w^{(n-1)/2}\).
Suppose the splitting character on both legs at stage \(j\) is
\(\chi_j\). Write \(B^{(0)}=B\), and at each unramified place form
the transferred label by
\[
 B_w^{(j+1)}=B_w^{(j)}\sqcup
    \{\chi_j(\varpi_w)q_w^{(h-1)/2-\ell}:0\le\ell<h\}.
\]
Here \(\sqcup\) adds multiplicities. This is the spherical transfer
in the normalized Satake coordinates, with the unramified character
twist restored \cite[Appendix, Proposition A.5 and the split formula
following it]{liu-spherical-core}. If the unitary characters \(\chi_j\)
have pairwise distinct infinity types, the spherical quotient in the
doubling integral at \(s=h/2\) is finite and nonzero.
\end{auxlemma}
\begin{proof}
For the initial spherical unitary representation the normalized
spherical coefficient is bounded by one. The bounded-spherical-function
criterion puts its real exponent in the convex hull of the Weyl
orbit of the half-sum of positive roots. In standard base-change
coordinates the extremal exponents at \(w\) are \(\pm(n-1)/2\).
This gives the assumed bound on \(B_w\) on the split as well as the
unramified unitary group. \cite[Section 4.7]{macdonald-bounded-spherical}

Write \(L^S(s,B)\) for the partial standard Euler product, with
unramified local factors \(\prod_{\alpha\in B_w}(1-\alpha q_w^{-s})^{-1}\).
Insert the displayed multiset in the numerator at \((h+1)/2\).
As \(\ell\) varies, the arguments of the new character factors run
through \(1,\ldots,h\). With \(\eta=\eta_{E/F}\) the quotient becomes
\[
 \frac{L^S((h+1)/2,B\otimes\chi_i^{-1})
    \displaystyle\prod_{j<i}\prod_{t=1}^{h} L_E^S(t,\chi_j\chi_i^{-1})}
 {\displaystyle\prod_{t=1}^{N_i} L_F^S(h+t,\eta^{h+t})}.
\]
The first product converges absolutely since
\((h+1)/2-(n-1)/2=3\). Absolute convergence handles the factors
with \(t>1\) and the denominator. At \(t=1\) the quotient character
is nontrivial. Hecke's nonvanishing theorem on the line
\(\Re s=1\) applies. Deleting finitely many unitary local factors
does not affect the conclusion.
\end{proof}

\begin{auxlemma}[From the norm of a lift to the span of its descents]\label{lem:fitted-core-span}
For the core chain of Lemma~\ref{lem:fitted-core-euler}, let
\(M_i(B')\) be the finite-dimensional multiplicity space of the
chosen real cuspidal type and finite level at label \(B'\).
Use the common split test of Lemma~\ref{lem:fitted-split-row-projector}
and the spherical vacuum away from the exceptional places.
Use the signed parameters of the core grid: in the centered
determinant-cover coordinates its upward pair is
\[
 (A,B\,;C,D),\qquad (A,S_h,D\,;C,B),\qquad
 S_h=((h-1)/2,(h-3)/2,\ldots,-(h-1)/2).
\]
Here \(A,C\) are positive, \(B,D\) negative, and the old coordinates
have absolute value at least \((h+3)/2\). The common splitting
character is included before centering. Use unramified vacua
at the unmarked finite places and the saturated level on the
specified split halves. Then the algebraic span of successive
descents from the full definite column space contains each bottom
row of the stated type, at its stated level.
\end{auxlemma}
\begin{proof}
\proofstep{First occurrence in the indicated towers.}
The old bottom entries and each previously inserted block lie outside
the two positions adjacent to the new block: the old and new centers
are separated by the grid gap. Thus \(S_h\) is the maximal consecutive
central string of the new parameter. With
\(r_*=|A|+|D|\) and \(s_*=|B|+|C|\), the signatures in the target
tower are
\[
       (r_*+h+l,s_*+l),\qquad l\in\mathbb Z .
\]

Paul's Theorem 2.7(b) identifies the signed lift
\cite{paul-real-core}. In the notation of Ichino's Theorem 6.1(ii)
the smaller centered parameter has \(k_\pi=0\), \(t=h/2>0\).
The sets \(C^\pm_\pi(t)\) are empty because each old entry lies
outside the interval of size \(t\). His condition \(l\ge k_\pi\)
holds at zero and fails below it:
\[
     \theta_{r_*+h,s_*}(\pi)\ne0,\qquad
     \theta_{r_*+h-1,s_*-1}(\pi)=0
\]
when the smaller signature exists. In the reverse direction deletion
of the maximal string is first occurrence by Paul's Proposition 3.4(ii).
The constant-term formula for theta lifts then gives cuspidality:
a nonzero constant term on a proper parabolic would give a lift to
a smaller member of the same tower. At the indicated real place its
local theta lift is zero. \cite[Theorem 6.1]{ichino-real-core}
\cite[Sections 1.2 and 11.4]{gqt-rallis}

\proofstep{The real and finite factors of the norm.}
For a matching tempered
type Ichino's Proposition 7.1 and Remark 7.2 give a nonzero local
doubling integral on Fock polynomials and compact-finite row vectors.
Polarizing first in the oscillator data and then in the row data gives
some \(P,u\) for which
\[
   \int_{U(C_i)(\mathbb R)}
       \langle \omega(g)P,P\rangle\,
       \langle \pi(g)u,u\rangle\,dg\ne0.
\]

Write \(D_i\) for descent on multiplicities with the finite test fixed.
The real Howe intertwining line factors the descent on real types as
\(D_i\otimes T_i(P)\). The Petersson adjoint on multiplicities is
\(D_i^*\), by interchanging the two absolutely convergent cusp integrals.
At the marked places the local coefficient is constant on the level
and zero off it. 
\proofstep{Apply the convergent doubling formula at every arrow.}
Let \(w_i\) be the Witt index of \(C_{i+1}\) over \(F\).
A rational totally isotropic subspace stays isotropic over every real
completion. Its dimension is at most the negative index at the selected
real place, hence
\[
     w_i\le m/2<h,\qquad N_{i+1}-w_i>N_i,\qquad
                           s_0=(N_{i+1}-N_i)/2=h/2 .
\]
Yamana's Theorem 2 on the split doubled space of dimension \(2N_i\)
identifies the theta integral with the Eisenstein series at \(s_0\)
\cite{yamana-siegel-weil}. Pair this identity with the two cusp forms on the two factors of
the doubled group. In the embedding with the second form negated,
the second factor carries the inverse determinant character
\(\chi_i^{-1}\circ\det\). Its cancellation in the diagonal theta kernel reads
\[
 (\chi_i^{-1}\circ\det)(g_2)(\chi_i\circ\det)(g_2)=1
                \qquad(g_2\in U(C_i)(\mathbb A_F)).
\]

 The product of the two theta
kernels on the diagonal gives the squared norm. Unfolding the
Eisenstein series first for large \(\Re s\), the boundary orbits pair
with constant terms of the cusp forms and vanish. The open orbit gives
the product of the local doubling integrals, by restriction of the
section to \(g\mapsto(g,1)\). This is the doubling computation in
\cite[Sections 11.3--11.5]{gqt-rallis}.

At a marked place the doubled height is one on the level. On its
complement the chosen local coefficient at \(s_0\) vanishes.
For the standard flat section the continuation there is explicitly
\[
       \Phi_s(g,1)=H(g)^{s-s_0}\Phi_{s_0}(g,1).
\]
Thus the marked support is independent of \(s\). At infinity write
\(H(g)\le1\) for the doubled Iwasawa height at \((g,1)\). 
The Siegel--Weil section at \(s_0\) on the doubled compact is bounded on
its finite-dimensional compact orbit. Its restriction for nearby
\(s\) satisfies
\[
        |\Phi_s(g,1)|\ll H(g)^{N_{i+1}/2+\Re s-s_0}
            =H(g)^{N_i/2+\Re s}\le1\quad(\Re s\ge0).
\]
The real discrete-series coefficient is integrable by the stipulated
noncompact-root gaps. Dominated convergence on compact substrips
therefore gives the local value at \(s_0\); at that point it is
Ichino's convergent matrix-coefficient integral. The unramified
quotient continues there by Lemma~\ref{lem:fitted-core-euler}.
This use of the convergent Siegel--Weil formula includes the longer
intermediate spherical strings.

On the fixed real vector the coefficient of an
occurrence \(a\) is \(\|a\|^2\langle\pi(g)u,u\rangle\) by the
definition of the multiplicity norm. The doubling identity evaluated below and Lemma~\ref{lem:fitted-core-euler} give
\[
 \|\Theta^\uparrow_{i,P}(a(u))\|^2
       =c_i(B')\,\|a\|^2,\qquad c_i(B')>0 .
\]

Indeed the spherical part is nonzero, the marked part is a product
of positive volumes and squared norms, and the real integrals are
nonzero. The norm identity gives the sign of their product. This
calculation applies to linear combinations in one multiplicity space:
averaging the two spherical arguments at a place acts by the same
Hecke scalar on every occurrence. In particular
\[
   \ker D_i^*|_{M_i(B')}=0,\qquad
   D_i^* M_i(B')\subset M_{i+1}(B'\sqcup S_{h,i}),
\]
where \(S_{h,i}\) denotes the new twisted string at all good places.

Removing the fixed new string with its multiplicities recovers the old label. Thus distinct initial labels
stay distinct. The composite adjoint is injective on their direct sum,
so in finite dimension \(D_0\cdots D_{m-1}\) maps onto the bottom
multiplicity space.

The span of \(T_i(P)v\), with \(P\) polynomial and \(v\) compact-finite,
is nonzero and stable under the lower maximal compact group.
Differentiating equivariance preserves this span because the
oscillator Lie algebra sends polynomials to polynomials. Irreducibility
makes it the whole lower Harish--Chandra module. Iterate this algebraic
span starting with a preimage under the composite multiplicity map.
At the definite end a map \(a:E_m^\vee\to\mathcal A(J_m)\)
corresponds in dual bases to the actual column
\[
                 g\longmapsto \sum_j a(e_j^\vee)(g)e_j .
\]
Each term in the resulting finite expansion is consequently a descent
from a permitted column.
\end{proof}

\begin{auxlemma}[Depth characters and the recovery of a lower test]\label{lem:fitted-lower-test-recovery}
In a polarized adelic Weil representation write \(T(v)\) for Weyl
translation by \(v\) in the symplectic space \(W(\mathbb A_F)\).
Let a lattice of finite translations and modulations \(L\) act by a
fixed character on a depth space. For any nonzero Schwartz function
\(f_0\) in this space, every Schwartz function \(f\) in the same
space is recovered by
\[
 f=\|f_0\|^{-2}\int_{L^\perp}
          \langle f,T(v)f_0\rangle\,T(v)f_0\,dv .
\]
The perpendicular denotes the symplectic character annihilator of
\(L\), with all real directions included. Measures on a polarization
and its dual are paired by Fourier inversion.
\end{auxlemma}
\begin{proof}
In position variables write \(v=(a,b)\), with \(a\) a translation
and \(b\) a modulation. The Weyl half-phases in the two factors cancel.
Fourier inversion in \(b\), followed by integration in \(a\), gives
at a position \(x\) the value
\[
              f(x)\int |f_0(x-a)|^2\,da=f(x)\|f_0\|^2 .
\]
The coefficient has compact support in the finite variables and is
Schwartz at infinity. Each Schwartz seminorm of \(T(v)f_0\) grows at
most polynomially at infinity, so the integral converges in the
Schwartz topology. Commuting \(T(v)\) past \(T(l)\), \(l\in L\),
multiplies the coefficient by \(\psi(\langle v,l\rangle)\) because
the depth characters of \(f\) and \(f_0\) agree. Hence it vanishes
outside \(L^\perp\).
\end{proof}

In particular, if the mixed Fourier restriction of an admissible
upper test is \(f_0\ne0\) and every \(T(v)f_0\) appearing here can
be realized by the same restriction at the prescribed upper level,
then any continuous lower period vanishing on the restricted
tests vanishes on the entire lower depth space. The support and
relative character are computed for each marked frame in Lemmas~\ref{lem:fitted-marked-support} and~\ref{lem:fitted-fixed-ascent}.

\begin{auxlemma}[Support of the marked Fourier coefficient]\label{lem:fitted-marked-support}
At an arrow write \(V=Y_{i+1}, W=Y_i\). In a Witt decomposition
\(V=X\oplus V_0\oplus Y\), the outer spaces are paired and
\(\ell=\dim_E X\). The hyperbolic pair \(X\oplus Y\) is orthogonal to \(V_0\). Let \(x\) be the ordered lower isotropic frame at
\(i=2j+1\); at \(i=2j+2\) append the distinguished nonisotropic
vector. Put
\[
 U=U_\ell^{\rm un},\qquad
 R_x=\{g\in U(W):gx=xu\text{ for some }u\in U\},
\]
where \(U_\ell^{\rm un}\) is the upper unitriangular group over \(E\).

For the finite level \(J\) of the lower row there are permissible
upper tests whose twice-transformed coefficient at the marked
arguments, integrated over \(u\) against the simple-root character,
is nonzero. For these tests this coefficient after translation by
\(g\) has support in
\[
 R_x(\mathbb A_F) U(W)(F_\infty) J.
\]
The same support holds after any real Weyl translation and after any
finite Weyl translation in the annihilator of the remaining depth lattice.
\end{auxlemma}
\begin{proof}
In the oriented homomorphism model let \(A_1\) be the projection
of the depth lattice to the first frame coordinate and \(A_2=A_1^\perp\)
its character annihilator in the paired coordinate. In the second
pair let \(B_1,B_2\) be the projections of the ordinary homomorphism
lattice with target shifted by \(\mathfrak j=\mathfrak m\bar{\mathfrak m}^{-1}\).
For the trace character and ideal conjugation, respectively,
\[
 B_2=B_1^\perp,\qquad \mathfrak j\bar{\mathfrak j}=\mathcal O_E.
\]
These identities give the two duality relations. Indeed in orthogonal matrix blocks
the ordinary homomorphism lattice is self-dual, and multiplication
of its target by \(\mathfrak j\) preserves this property. Write
\(y=-ze_\ell^t\) for the adjoint argument in the odd step and \(y=0\)
in the even step; empty frame coordinates are omitted.

\proofstep{Conditions forced by a nonzero local value.}
In the two Fourier coordinates use finite factors
\(\mathbf1_{x+A_1}\) and \(\mathbf1_{y+B_2}\). Translation and
modulation in these coordinates give the following necessary conditions
at each finite place:
\begin{align*}
 g^{-1}xu-x&\in A_1\quad\text{for some }u\in U,\\
 g^{-1}xN\in A_1&\ \Longrightarrow \psi_E\Big(\sum_{j<\ell}N_{j,j+1}\Big)=1
 \quad (N\text{ strictly upper triangular}).
\end{align*}
The first follows by pairing with all of \(A_2\). For the second,
the substitution \(u\mapsto u+N\) preserves the product Haar measure
on the upper entries. In the second pair, whenever \(g\xi=\xi'\)
are vectors in the same isotropic coordinate and \(\xi\in B_1\),
modulation also forces \(\psi(\langle y,\xi'-\xi\rangle)=1\).
These conditions hold on \(gJ\), since the finite tests are fixed
by the lower level.

\proofstep{Selected split places.}
Write \(I=\varpi^d\mathcal O\) for the selected half of
\(\mathfrak m\). Denote the ordinary coefficient ideal by
\(\mathfrak A\), so an ordinary map sends
\(\Lambda_V\) into \(\mathfrak A\Lambda_W\), and put
\[
 \mathfrak a=\mathfrak A\bar{\mathfrak m}^{-1},\qquad
 \bar{\mathfrak a}=\bar{\mathfrak A}\mathfrak m^{-1}.
\]
Both halves of \(\mathfrak A\) are units at a selected place. In a
block of size \(c\), column \(j=0,\ldots,c-1\) has coefficient ideal
\(\mathfrak a^{c-j}\bar{\mathfrak a}^{-(j+1)}\).
Its frame vector carries the scalar \(\eta_j=(-\xi_0)^{j+1}\).

The upper block has one more column. In integral bases the marked
vectors and the covectors on the opposite half are therefore
primitive. The point condition gives congruence modulo \(I\) on
both halves: on the selected half the map itself has depth \(I\);
on the opposite half the two column ideals have ratio \(\varpi^d\),
since
\[
  \varpi^{-d(c-j)}
        =\varpi^d\varpi^{-d(c+1-j)} .
\]

Absorb the antidiagonal pairing signs in the covector bases. If
\(A,B\) and \(A',B'\) are the original and transformed inclusions
and surjections, then
\[
 A'\equiv A,\quad B'\equiv B\pmod I,\qquad B'A'=BA.
\]
The Gram equality holds also for the appended line since the last
row of \(u\) is fixed. Choose \(L A=1\), \(B'J'=1\). The corrections
\[
 T_0=1+(A'-A)L,\qquad T=T_0+J'(B-B'T_0)
\]
give \(TA=A'\), \(B'T=B\) and \(T\equiv1\pmod I\).
The split unitary transformation defined by \(T\) and its inverse
dual therefore puts \(g\) in \(R_x J_v\).

\proofstep{Places away from the depth.}
Let \(J_j\) be the coefficient ideal of the marked \(j\)-th upper
frame vector, so that scalar multiplication by \(J_j\) makes it
primitive in the upper lattice. Projection to each Witt coordinate
preserves the lattices. Off the depth the successive ideal ratio is
\[
 J_j J_{j+1}^{-1}=(\xi_0)\mathfrak A\bar{\mathfrak A}
                    =\operatorname{Ann}_{\psi_E}(\mathcal O_{E,v}).
\]
This follows by dividing the preceding column ideals and multiplying
by \(\eta_{j+1}/\eta_j\). Use the Witt parabolic decomposition
\(g=n d(t)\gamma k\), with unipotent \(n\), diagonal parameters
\(t_j\), core isometry \(\gamma\), and full lattice stabilizer
\(k\). Integrally, this decomposition follows by extending a
primitive isotropic flag to a hyperbolic basis.

The grid field extension is unramified at every finite place and split above 2. At a split place take dual bases on the two halves. Else choose vectors of Gram matrix \(H\) dual to the given frame. Subtracting half of \(H\) times the frame gives
\[
 H-\tfrac12 H-\tfrac12 H^*=0.
\]
The resulting dual vectors are isotropic.

In the point condition project to the \(j\)-th isotropic line.
The earlier columns vanish under projection. Primitivity gives
\(t_j^{-1}J_j\subset J_j\). In the shift condition use the simple
matrix \(t_j s E_{j,j+1}\), multiplied on the left by the
unitriangular matrix induced by \(n\). For \(s\in J_j J_{j+1}^{-1}\)
its translated image is an integral map, and the sum of adjacent
entries is \(t_j s\). Thus
\[
 t_j^{-1}\in\mathcal O_{E,v},\qquad
 \psi_E(t_j J_jJ_{j+1}^{-1})=1
           \ \Longrightarrow\ t_j\in\mathcal O_{E,v}.
\]

This reaches every hyperbolic column in the even case. In the odd
case one last column remains when its isotropic rank \(q\) is positive.
Its ideal and the ideal of the upper distinguished line are
\[
 \mathfrak b=(-\xi_0)^{-q}\mathfrak A^{q+1}\bar{\mathfrak A}^{-q},
 \qquad
 \mathfrak c=(-\xi_0)^{-(q+1)}
           \mathfrak A^{q+1}\bar{\mathfrak A}^{-(q+1)}.
\]

Dividing these expressions gives
\[
 \bar{\mathfrak A}\mathfrak b^{-1}=(\xi_0^{-1})\mathfrak c^{-1}.
\]
 The covector of the line character,
scaled by \(s\) before adjunction, has adjoint image ideal
\((\xi_0^{-1})\bar s\bar{\mathfrak c}^{-1}\);
the target dual line has ideal \(\bar{\mathfrak b}^{-1}\) before
ordinary homomorphism scaling. By the displayed identity this
adjoint is integral precisely for \(s\in\mathcal O_{E,v}\).
Apply the second-block character condition after removing \(k\)
by right invariance. On the last column the unipotent has diagonal
one. The scaled covector for \(s=t_q-1\) annihilates the entire
complementary coordinate, hence lies in its dual lattice. Therefore
\[
        t_q-1\in\mathcal O_{E,v},\qquad t_q\in\mathcal O_{E,v}^{\times}.
\]

In the even case projection to the remaining Witt core gives
\(L_z\subset\gamma\Lambda_0\) for the distinguished self-dual line
lattice \(L_z\). Orthogonal projection preserves the unimodular
lattice \(\gamma\Lambda_0\), so
\[
 \gamma\Lambda_0=L_z\perp M,\qquad \Lambda_0=L_z\perp M_0
\]
with \(M,M_0\) self-dual in the same complement. Over an unramified
quadratic extension of odd residue characteristic the residual
nondegenerate Hermitian form has a vector of nonzero norm. Split
it off, use surjectivity of the norm on units, and repeat to obtain
an orthonormal integral basis. In the split case the second lattice
is dual to the first half. An isometry of the complements extending
the identity on \(L_z\) gives \(\gamma\in R_xJ_v\). In the isotropic
case the whole core is already in \(R_x\). Consequently all local
factors lie in \(R_xJ_v\), with \(R_x=U(W)\) for the empty frame \(\ell=0\).

Since \(g_v\in J_v\) at almost all places, the local decompositions give
\[
 g_f=(r_v)_v(k_v)_v,\qquad (r_v)_v\in R_x(\mathbb A_{F,f}),\quad(k_v)_v\in J.
\]


\proofstep{Construction at the marked arguments.}
The two hyperbolic pairs and the residual block split the ordinary
and depth lattices by the integral coordinate projections. The
marked frame belongs to the trace annihilator of the first dual
depth component. On the selected pair its two bounds are the
ordinary lattice and its inverse-depth enlargement. For the second
argument at a selected place, retaining the modulus in the computation
of \(\mathfrak b,\mathfrak c\), with \(\mathfrak a,\bar{\mathfrak a}\)
in place of \(\mathfrak A,\bar{\mathfrak A}\), gives
\[
       \bar{\mathfrak A}\mathfrak b^{-1}\mathfrak m^{-1}
                         =(\xi_0^{-1})\mathfrak c^{-1}.
\]
Thus both marked eigencharacters agree with the depth action.
At a dyadic split place impose the translation and modulation
conditions on their separate halves. In Schrödinger coordinates
they give the span of the functions \(\mathbf1_{t+I X}\) on \(X\);
off the depth use the self-dual vacuum.

In the two mixed coordinates take the preceding finite indicators,
real bumps positive at their marked arguments, and a nonzero
residual depth vector \(f\). Inverse Fourier transform in the two
outer pairs gives an upper Schwartz test. The coefficient it yields
is the product of the second bump at \(y\), the vector \(f\), and
\[
       \int_U b_1(xu)\,
             \overline{\psi_E(\textstyle\sum u_{j,j+1})}\,du .
\]
The finite phase is one on the support: the ratio of the lower
\(j\)-th ideal to the next upper one is \((\xi_0)\bar{\mathfrak a}\).
Multiplication by the scaled target ideal, including the one-sided
depth where selected, puts each adjacent entry in
\((\xi_0)\mathfrak A\bar{\mathfrak A}\).

The map \(u\mapsto xu\) is proper in the upper entries by a left
inverse of \(x\). Choose the real bumps sufficiently small to make
the real part of the phase positive. The integral is nonzero.
The constant depth characters are finite; real Weyl translations
preserve them, and residual translations by the depth annihilator
commute with both outer Fourier transforms.
\end{proof}

\begin{auxlemma}[Ascent preserving a specified descending chain]\label{lem:fitted-fixed-ascent}
Use the auxiliary grid of Definition~\ref{hyp:fitted-grid-data}.
Fix a sequence of cuspidal core forms \(c_i\) with actual downward
tests \(D_i\), and a fixed smooth bottom row \(G_0\), satisfying
\[
 c_i=\Theta_{D_i}(c_{i+1}),\quad 0\le i<m,\qquad
 \int_{[U(Y_0)]} c_0(\iota h)G_0(h)\,dh\ne0 .
\]
Here \([L]=L(F)\backslash L(\mathbb A_F)\) for an \(F\)-group \(L\).
Keeping \(c_m\) and the bottom vector fixed, some allowed finite coset and real Fock polynomial
at each arrow give a nonzero terminal Bessel period.
\end{auxlemma}
\begin{proof}
Use the conjugate realization, reversing the sign of the first
Hermitian form; the resulting additive character is the original
\(\psi\). Set \(\psi_E=\psi\circ\operatorname{Tr}_{E/F}\). Write
\(\alpha_i\) for the inverse splitting character of grid index \(m-i\)
and \(\eta\) for the odd line splitting. In the orientation with the
larger space first, the two row splitting characters are
\[
 (\alpha_i\eta,\alpha_i)\quad(i\text{ odd}),\qquad
 (\alpha_i,\alpha_i\eta)\quad(i\text{ even}).
\]
On the core they are \((\alpha_i,\alpha_i)\). Orthogonal direct sum
multiplies the two characters on the side being decomposed. The
remaining line thus has precisely the stated splitting \(\eta\).
The adjoint map to the linked line reverses the symplectic form
and transports the action by complex conjugation.

\proofstep{Cuspidality of the intermediate row lifts.}
Write \(s=r+6\) for the odd increase in row rank and
\(\delta_i=-1/2\) for odd \(i\), \(\delta_i=0\) for even \(i\).
Let \(P_i\) be the unshifted signed row list, so that the grid uses
the dual of \(P_i+\delta_i\). On the arrow with new center \(c\),
the splitting exponents on the two dimensions are
\(c+\delta_i,c+\delta_{i+1}\). In centered coordinates all old
entries are below the new string:
\[
       x+(s+1)/2\le c\quad(x\in P_i),\qquad
       S_s=((s-1)/2,(s-3)/2,\ldots,-(s-1)/2).
\]
Indeed the new center lies above the old lists by the separating gap.

Negating the larger form exchanges the old signature labels and places
\(S_s\) in its negative list. Paul's Theorem 2.7(b) gives this signed
pair on the determinant covers, and taking duals restores the chosen
row. In Ichino's notation there is no central zero, and
\[
 k_\pi=-1,\qquad t=(s-1)/2,\qquad C^+_\pi(t)=C^-_\pi(t)=\varnothing .
\]
 Thus the local condition
in this tower is satisfied at \(l=0\) and fails at \(l=-1\). The first
occurrence constant-term formula gives the asserted cuspidality of
every nonzero partial row lift. \cite{paul-real-core}
\cite[Theorem 6.1(i)]{ichino-real-core}

\proofstep{The rational orbit in each Fourier coefficient.}
Use the notation \(X,V_0,Y,x,U,R_x\) of
Lemma~\ref{lem:fitted-marked-support} and put
\(\chi(u)=\psi_E(\sum_{j<\ell}u_{j,j+1})\).
Write a unipotent element as \(n_C(c)n_B(b)m(u)\), where
\(c^*=-c\), \(b:V_0\to X\), and \(u\in U\).
In the lattice realization the Weyl operators satisfy
\[
 T_s f(t)=\psi(B(s,t)/2)f(t+s),\qquad
 T_sT_v=\psi(-B(s,v)/2)T_{s+v}.
\]
Here \(B\) is the symplectic form on the homomorphism space.
Let \(F^X(x';t)\) be the first mixed coefficient of an oscillator
test \(F\): Fourier expand its evaluation at zero in the compact
quotient of the dual \(Y\)-coordinate, take the coefficient at
\(x'\), and leave the middle variable \(t\) free.

The central \(c\)-integration selects Gram matrix zero for an upper
Bessel period, and \(\operatorname{diag}(0,\ldots,0,\langle z,z\rangle)\)
for an upper Fourier--Jacobi period. Indeed its phase at \(x'\)
is \(\psi(B(j_Xx',j_Y(x'c))/2)\), whose linear dependence on
\(c\) records the Gram matrix by the nondegenerate trace pairing.
The two resulting Gram constraints are
\[
 (x')^*h_W x'=0\quad\text{or}\quad (x')^*h_W x'=\langle z,z\rangle e_\ell e_\ell^t.
\]
The second comes from the conjugate line phase at \(ze_\ell^t\).

If rational columns are dependent, choose the first relation
\(x'd=0\) with \(d_j=1\), \(d_k=0\) for \(k>j\). When \(j<\ell\)
the stabilizer contains \(1+dse_{j+1}^t\), fixing the frame and
last row but contributing \(\psi_E(s)\). When \(j=\ell\), the
nonzero last-Gram case is impossible. In the other case take the
side translation \(b=sdz^*/\langle z,z\rangle\), giving
\[
 x'b=0,\qquad \int_{E\backslash\mathbb A_E}\psi_E(s)\,ds=0 .
\]
Witt extension makes the independent frames with the specified
Gram matrix a single rational orbit. After inserting its
representative \(x\), the stabilizer modulo changes of upper
unitriangular basis is \(R_x\).

\proofstep{Two different contractions of the surviving orbit.}
Write \(G=G_i,\ c=c_{i+1}\). For upper core coordinate \(a\), lower
group coordinate \(g\), and upper oscillator test \(\Phi\), put
\(F_{a,g}=\omega(\iota(a),g)\Phi\). Unfolding the rational orbit
writes the upper period as
\begin{align*}
 &\int_{R_x(F)\backslash U(W)(\mathbb A_F)} f_\Phi(g)\,dg,\\
 & f_\Phi(g)=G(g)\int_{[U(C_{i+1})]}c(a)
  \int_{U(\mathbb A_E)}\int_{[b]}\theta_a(b,u)
          F_{a,g}^X(xu;xb)\,db\,du\,da .
\end{align*}
The notation \([b]\) denotes the rational quotient of the additive
side block. Unipotent quotients have mass one. Positive rescaling
of the quotient measures multiplies the corresponding period by
a positive constant.

For upper Bessel, set \(y=-ze_\ell^t\) in the adjoint block.
The weight is \(\overline{\chi(u)\psi_E(e_\ell^t b z)}\).
The lattice relation at rational \(y\) and the two Weyl half phases
give, for the embeddings \(j(b)\) of the side and \(j^*(y)\) of
its adjoint,
\[
 \overline{\psi_E(e_\ell^t b z)}F^X(xu;xb)
      =(T_{j(b)}T_{j^*(y)}F^X(xu;\cdot))(0).
\]

Integration in \([b]\) is consequently the second mixed coefficient
at \(y\); integration in \(u\) uses \(\bar\chi\). Let \(v=j^*(y)\).
For lower flag element \(n\), substitution of \(nx=xu_n\) gives
the simple-root factor from \(\chi(u_n)\); the second argument becomes
\(n^{-1}v\). Moving its rational part past evaluation leaves the
scalar \(\psi(B(v,n^{-1}v)/2)\) and translation by the line component
of \(n^{-1}v-v\). These are exactly the central and side actions
of the lower Jacobi line. The remaining two middle groups commute
with the outer evaluations and act by the core Weil representation.
Write \(\mathcal M\) for this mean and \(\mathcal P_{\rm FJ}\) for the lower Jacobi period. The calculation gives
\[
 \mathcal M(D_i\otimes L)=\mathcal P_{\rm FJ}(G_i,c_i;L),\qquad c_i=\Theta_{D_i}(c_{i+1}).
\]



For upper Fourier--Jacobi let \(Y_a=\omega_{\rm line}(a,1)\Phi_0\),
where \(\Phi_0\ne0\) is the fixed upper line test. Here the weight
is \(\overline{\chi(u)Y_a(ze_\ell^t b)}\). The first \(\ell-1\)
side rows extract the second mixed coefficient at zero. The last
row maps isomorphically onto the rational line block, preserving
its mass-one quotient measure. If \(\Xi\) denotes the resulting
test before contraction and \(h\) belongs to the lower middle group,
the remaining integral is
\[
 \int_{[\mathcal W_{\rm line}]}
     (\omega_0(a,(h,1))\Xi)(0,t)\,\overline{Y_a(t)}\,dt .
\]
For the appended frame the last adjacent entry of \(u_n\) is the
lower side entry \(e^t b z\). The other adjacent entries give the
lower flag character.

Thus the \([R_x]\)-mean is the lower Bessel
period. At \(\Xi=D_i\otimes\Phi_0\) the line integral gives
\(\|\Phi_0\|^2\) by unitarity, and integration in \(a\) gives the lower form. With \(\mathcal P_{\rm B}\) the lower Bessel period,
\[
 \mathcal M(D_i\otimes\Phi_0)=\|\Phi_0\|^2\mathcal P_{\rm B}(G_i,c_i).
\]
The norm is positive because the same nonzero vector occurs in the two line slots.

\proofstep{Passage from the global period to its lower mean.}
These integrations are absolute for cuspidal smooth input. In a
fixed finite level and support, the real lattice sum of a translated
Schwartz function is bounded using
\(\sum_{k\in\mathbb Z^N}(1+\|k\|)^{-N-1}<\infty\).
After adelic translation by \(g\) the index of a denominator lattice
for this sum is bounded by a power of matrix height: intersect the
translated lattice with a fixed one and bound denominators by the
entries of the inverse action. Real Cartan decomposition bounds
every needed translated Schwartz seminorm polynomially in height.
The compact unipotent integrations preserve these estimates.

In particular, for the kernel sum \(S_\Phi(a,g)\) before unfolding, some Schwartz seminorm \(p\) and fixed exponents give
\[
 |S_\Phi(a,g)|\ll p(\Phi)\operatorname{ht}(a)^A
                               \operatorname{ht}(g)^B .
\]

For the restricted mean, let \(e_X\) be the wedge of the rational
isotropic basis. On \(e_X\wedge V_0\) the unipotent is trivial
and the middle element \(h\) acts by \(e_X\wedge w\mapsto
e_X\wedge hw\). For \(w\in V_0(F)\setminus\{0\}\), the product
formula applied to any rational translate of \(e_X\wedge w\) bounds
\(\|h^{-1}w\|\) below by an inverse power of ambient quotient height.
Reduction of the middle lattice therefore bounds its quotient height
by a power of the ambient one. For the subgroup fixing a line use
the remaining direct summand. The reduced heights of the middle and ambient groups therefore satisfy
\[
 \operatorname{ht}_{\rm mid}(h)\ll(1+\operatorname{ht}_{\rm amb}(n\iota(h)))^{A_1}.
\]
Rapid decay on both reductive quotients dominates these powers. These estimates also
give an integrable majorant on compact sets of right translates
and continuity in the real Schwartz test.

Suppose every permitted upper test gives zero. Real Hermite functions
span densely in the Schwartz topology; finite coset functions span
the selected depth space. This vanishing therefore holds after real
Weyl translations of the constructed tests. Modulating the first
dual coordinate at infinity multiplies the unfolded integrand by
\[
       \psi_\infty\big(B((g^{-1}xu)_\infty,t)\big)
       \qquad(t\text{ any real dual coordinate}).
\]
The pushforward under \((g,u)\mapsto(g^{-1}xu)_\infty\) is thus a finite
complex measure with identically zero Fourier transform. For any
compact \(Q\subset U(W)(F_\infty)\), the set \(Q^{-1}xU(F_\infty)\)
is closed: pass to a convergent subsequence in \(Q\) and use a left
inverse of \(x\) on the upper entries. Its preimage imposes exactly
\(g_\infty\in R_x(F_\infty)Q\). Every such strip has zero integral.

By Lemma~\ref{lem:fitted-marked-support} only one finite double
coset remains and the integrand is right \(J\)-invariant.
The unipotent radical and the unitary middle group of \(R_x\) are
unimodular; the modulus on the radical is a rational character of
the middle group and is trivial. Disintegration on the support
multiplies the real integral of the \([R_x]\)-mean by a positive
finite factor proportional to
\[
   \operatorname{vol}(J)/
   \operatorname{vol}(J\cap R_x(\mathbb A_{F,f})).
\]
The mean is continuous, so its value at the identity vanishes.
If the frame rank is zero the real frame space is a point and
the same statement follows directly by quotient integration.

\proofstep{The two prescribed lower tests.}
The residual lattice separates the core and line projections.
At a stage indexed by \(2\rho\) let
\(\mathfrak l_\rho=\mathfrak a^\rho\bar{\mathfrak a}^{-\rho}\) and
\(\mathfrak j=\mathfrak m\bar{\mathfrak m}^{-1}\) as in
Lemma~\ref{lem:fitted-marked-support}. 
The underlying line lattice
has coefficient \(\mathfrak l_\rho\) and its test uses the target
\(\mathfrak j\mathfrak l_\rho\). Both are self-dual. At an even
arrow projection gives target \(\mathfrak j\mathfrak l_{\rho+1}\)
in the line factor. At an odd arrow the inverse line link is
\(\xi_0^{-1}\) times Hermitian adjoint, since
\(\zeta_{\rho+1}=(-\xi_0)\zeta_\rho\). The coefficient changes
from \(\mathfrak A\) to \(\bar{\mathfrak A}\), and
\[
    \bar{\mathfrak A}\mathfrak l_{\rho+1}
                =\mathfrak A\mathfrak j\mathfrak l_\rho .
\]

Thus the linked conjugate line test has the precise target at
positive \(\rho\). For \(\rho=0\), the ordinary and the shifted
upper targets impose respectively the selected-half depth and the
ordinary line condition on its adjoint. Together they give the
projected rank-zero line lattice. Every core projection is onto
an original orthogonal summand and respects its ordinary and
depth conditions. Weyl operators factor under the orthogonal
decomposition, so for their respective depth classes one has
\[
 D_i\otimes L\in\mathcal S_{\rm BF}^{\rm res},\qquad D_i\otimes\Phi_0\in\mathcal S_{\rm FB}^{\rm res}.
\]


Residual Weyl translations by the annihilator extend on the upper
test by acting in the core block; they preserve its two outer
characters. Lemma~\ref{lem:fitted-marked-support} gives a nonzero
residual vector in the image. Lemma~\ref{lem:fitted-lower-test-recovery}
therefore extends vanishing of the mean to the two specified
tensors. The contractions just computed would give zero for the
preceding period, which proves the contrapositive at each arrow.

At rank zero \(C_0=Y_0\perp L_z\), \(C_1=Y_1\).
Expand the bottom test in product tests on the two orthogonal blocks.
Their finite cosets span and their real Hermite sums converge
in Schwartz space. On a product, Fubini for the two theta kernels
identifies the bottom pairing with the first Fourier--Jacobi
period:
\[
 \int_{[U(Y_1)]} c_1(a)\theta_L(a,1)
       \left(\int_{[U(Y_0)]}\theta_D(a,h)G_0(h)\,dh\right) da .
\]

Some product has nonzero value. Repeated application of the two
contrapositives preserves the tests already chosen; density and
linearity always yield some individual finite coset and real
polynomial. At the last even stage the frame and line satisfy
\[
 e_j=(-\xi_0)^j e^{\rm grid}_j,\qquad z=(-\xi_0)^{\rho+1}z^{\rm grid}.
\]

Every adjacent coefficient, including the line entry, is multiplied
by \(-\xi_0\). Consequently the auxiliary simple-root character
is the conjugate of \(\psi_E(\xi_0\sum s_j)\); conjugation recovers
the prescribed terminal Bessel period, up to a positive Haar ratio.
\end{proof}

\begin{lemma}[A criterion using the prescribed test spaces]\label{lem:fitted-grid-ascent-criterion}
Fix the data and the row \(w\) of Definition~\ref{hyp:fitted-pointwise-nonvanishing}. The following two properties suffice for its pointwise assertion.
\begin{enumerate}
\item Write \(\mathcal C\) for the admissible bottom forms and \(B_0(F,v)\) for their period against the row of \((v,w)\). Require
\[
 \forall v\in V\setminus\{0\}\quad\exists F\in\operatorname{span}_{\mathbb C}\mathcal C
                  \qquad B_0(F,v)\ne0 .
\]
\item A nonzero bottom period gives a nonzero functional on unrestricted oscillator tests at the first Fourier--Jacobi step. At each subsequent Bessel or Fourier--Jacobi step, a nonzero preceding pairing likewise gives a nonzero functional on the next tests. Each functional is invariant under the prescribed compact open level in the finite variables, and its projected test space is spanned, modulo its kernel, by the stipulated finite tests and admissible real Fock polynomials.
\end{enumerate}
All the indicated pairings must converge on these test spaces.
\end{lemma}

\begin{proof}
If \(F=\sum_i c_iF_i\), linearity of the convergent bottom period gives
\[
 0\ne B_0(F,v)=\sum_i c_i B_0(F_i,v).
\]
At least one \(F_i\) therefore supplies an actual admissible column. At one transition let \(\Lambda\) be the nonzero functional, choose \(f\) with \(\Lambda(f)\ne0\), and let \(K\) be the compact open level. With probability Haar measure,
\[
  \Lambda\!\left(\int_K kf\,dk\right)=\int_K\Lambda(kf)\,dk=\Lambda(f).
\]
Smoothness makes the finite-adelic average a finite sum. Expand the projection modulo \(\ker\Lambda\) in the specified tests. At least one summand has nonzero value. Repeat first at the bottom step, then alternately at the Bessel and Fourier--Jacobi steps until the last even level \(m\). Its nonzero value is \(I_v(\rho,\xi;\Psi,\Phi)\) for the tests retained along this chain.
\end{proof}

\begin{auxlemma}[Pointwise nonvanishing on the prescribed grid]\label{lem:fitted-completed-nonvanishing}
Use the field, characters, principal level, signed parameters and
compatible measures of Definition~\ref{hyp:fitted-pointwise-nonvanishing}.
For \(C_F\ge6510\), above the scale threshold of
Lemma~\ref{lem:fitted-depth-allowance}, fix a nonzero decomposable
smooth \(w\). For each nonzero bottom multiplicity \(v\) there are
permitted tests at this very level with
\[
                    I_v(\rho,\xi;\Psi,\Phi)\ne0 .
\]
\end{auxlemma}
\begin{proof}
Take the dual real row for the bilinear period. At a real place merge
the two signed parameter lists in decreasing order:
\[
 a_0,\ldots,a_m,\ v_2,v_2-\tfrac12,\ldots,v_k,v_k-\tfrac12,          -b_1,\ldots,-b_m .
\]
The first block is in the core and alternates from positive. The last
is in the row dual and alternates from positive. The middle block
alternates row and core and has positive signatures throughout. Since
\(m\) is even, at both joins and within each block adjacent entries
have equal signs precisely when they belong to different lists.
The first and last entries have the signs for deleting a positive
line. He's signed interlacing theorem therefore gives the required
Hilbert summand with multiplicity one. \cite[Theorem 0.1]{he-discrete-branch}

The entries in the first ambient block satisfy
\[
        a_i=v_2+\tfrac12+(m+1-i)g_a\qquad(0\le i\le m).
\]
The middle block uses \(2\le j\le k\); the final one has
\(1\le l\le m\). At the adjacent indices in these ranges the grid inequalities give
\[
 v_j-v_{j+1}\ge2,\quad v_k+b_1\ge2,\quad
 \min(g_a,b_{l+1}-b_l)\ge1000n^3,\quad
                   \max(a_0,b_m)\le b^2.
\]
Within the first ambient block the spacing is \(g_a\), and between
its smallest entry and the second block the gap is \(1+g_a\).
The second block has constant signature. In the row-dual list
the first possible opposite-sign final entry after the middle block
is \(-b_2\); its gap is at least \(v_k+b_1+(b_2-b_1)\).
All opposite-sign pairs in the two outer blocks have at least the
indicated large spacing. Since \(0\le v_k\) and \(0<b_1\), the
coordinate bounds and the opposite-sign separation imply
\[
   \lambda_i,\mu_j\in[-b^2,b^2],\qquad
   2b^2\ge1000n^3,\qquad b\ge n .
\]

For the resulting row vector, Lemmas~\ref{lem:weighted-fitted-translate}
and~\ref{lem:fitted-real-projection} give the reproducing and derivative
estimates with one scalar. Lemmas~\ref{lem:fitted-large-row-lattice}
and~\ref{lem:fitted-prescribed-tail} then give a bottom Poincar\'e
series with nonzero period. Its real seed is a translate of an
integrable discrete-series coefficient: the two real translations
may be applied to the fixed lowest-type seed. The usual cuspidal
Poincar\'e construction in Lemma~\ref{lem:bottom-unbounded-depth}
applies at this finite level. Irreducibility, right translation and
compact averaging, as in the last step of that lemma, give a
compact-finite test pairing nontrivially with the specified \(w\).

For this fixed \(w\), put \(G_0=v(w)\). By Lemma~\ref{lem:fitted-core-span} write the core test as a finite sum of descents \(c_{0,j}\). If \(B_0\) is the bottom period, then
\[
  0\ne B_0(\textstyle\sum_j c_{0,j},G_0)=\sum_j B_0(c_{0,j},G_0).
\]
Choose a summand with nonzero period.
Lemma~\ref{lem:fitted-fixed-ascent} carries that very column and the
bottom row through the alternating steps. Its nonzero terminal
integral is the asserted value.
\end{proof}

\begin{construction}
\[
\begin{array}{ccccc}
 c_0 &\xleftarrow{\ \Theta_{D_0}\ }& c_1
     &\xleftarrow{\quad\cdots\quad}& c_m\\
 \big\downarrow&&\big\downarrow&&\big\downarrow\\[-2pt]
 B_0(c_0,G_0)\ne0 &\Longrightarrow&
      B_1(c_1,G_1)\ne0 &\Longrightarrow& B_m(c_m,G_m)\ne0
\end{array}
\]
\caption{The top row constructs a test for the original period from
a definite column. Along the bottom row each implication chooses
one permitted test for the next lift; the previously chosen column
and the bottom row vector stay fixed.}
\label{fig:fitted-period-chain}
\end{construction}

\begin{proposition}[A finite separating family at the base tuple]\label{prop:fitted-grid-base-separation}
For the data of Definition~\ref{hyp:fitted-pointwise-nonvanishing}, one can fix a single pair \((\rho_0,\xi_0)\) and choose \(d=\dim_{\mathbb C}V\) admissible tests \((\Psi_i,\Phi_i)\) for which
\[
   V\longrightarrow\mathbb C^d,\qquad
   v\longmapsto\bigl(I_v(\rho_0,\xi_0;\Psi_i,\Phi_i)\bigr)_{i=1}^d
\]
is an isomorphism.
\end{proposition}

\begin{proof}
At a terminal real place \(\tau\) fix all other variables in the Bessel integral. Its restriction in the two local representation variables is a continuous bilinear functional transforming by the prescribed generic Bessel character. Indeed, right translation by an element of the Bessel subgroup in both arguments is removed by a change of variables, leaving the character of the unipotent. The representations are the smooth discrete series and the irreducible type of the compact unitary companion. For this unitary Bessel pair write \(B_\tau\) for the Bessel subgroup and \(\xi_\tau\) for its prescribed character. Archimedean multiplicity one gives
\[
 \dim\operatorname{Hom}_{B_\tau}(\pi_\tau\widehat\otimes\sigma_\tau,\xi_\tau)\le 1
\]
for these irreducible smooth representations. \cite[Theorem~A]{jiang-sun-zhu-bessel}

If a terminal integral is nonzero, choose a generator \(b_\tau\) of its local functional line and a pair of vectors at which \(b_\tau=1\), at every real place. For pure tensors of local vectors \(z_\tau\), evaluate one place at a time at its chosen pair. Proportionality of the local forms gives, for each terminal integral \(B\),
\[
 B((z_\tau)_\tau)=B((z_\tau^0)_\tau)\prod_\tau b_\tau(z_\tau).
\]
Pure tensors span the coefficient spaces, so multilinearity extends the identity to sums.

 The real intertwiners of the fixed tower are common to its finite tests, as in \eqref{eq:down-theta}. Pulling the terminal real functional back by them and evaluating at the fixed \(w\) therefore gives a common function \(a_w\) with
\[
   I_v(\rho,\xi;\Psi,\Phi)=a_w(\rho,\xi)\,c(v;\Psi,\Phi).
\]
The case \(V=0\) uses the empty list. Otherwise apply Lemma~\ref{lem:fitted-completed-nonvanishing} to one nonzero vector to obtain \(a_w(\rho_0,\xi_0)\ne0\). For every other nonzero \(v\), a nonzero integral at any real test implies \(c(v;\Psi,\Phi)\ne0\). The same finite test thus works at \((\rho_0,\xi_0)\).

Starting with \(V_0=V\), choose a nonzero vector in \(V_i\) and a test functional \(L_{i+1}\) nonzero on it. Then
\[
 V_{i+1}=V_i\cap\ker L_{i+1},\qquad \dim V_{i+1}=\dim V_i-1 .
\]
After \(d\) choices the coordinate map is injective between spaces of equal dimension. This also shows exactly where the pointwise automorphic input enters the row-height estimate of Proposition~\ref{prop:descended-row-height}.
\end{proof}

\begin{auxlemma}[A determinant test on the realized base row]\label{lem:fitted-base-minor}
Let \(P:V_{t_0}\twoheadrightarrow R\) be the column map to the realized base row in Section~\ref{subsec:fitted-row-height}, and put \(s=\dim_KR>0\). Choose a finite list of admissible tests at a common real pair, with normalized base matrix \(M=V(t_0)\) from Proposition~\ref{prop:fitted-grid-coefficient-matrix}. Let \(F\) be their actual period map on \(R_\mathbb C\). Suppose the base comparison for this list is
\[
                  zc\,F P_{\mathbb C}=c_*\sigma_0(M),\qquad zc\ne0 .
\]
One nonzero value of \(FP_{\mathbb C}\) and one nonzero \(s\)-rowed minor of \(M\) suffice for these tests to separate every vector in \(R_\mathbb C\). For a common integral denominator \(a\ne0\) of \(M\), the minor condition is equivalent to
\[
     \operatorname{rank}_{\kappa(\mathfrak p)}(aM\bmod\mathfrak p)\ge s
     \quad\text{at some finite prime }\mathfrak p\text{ of }K .
\]
\end{auxlemma}
\begin{proof}
The nonzero period forces \(c_*\ne0\). Thus \(\operatorname{rank}_K M=\operatorname{rank}_{\mathbb C}F\) by surjectivity of \(P\) and injectivity of the field embedding. This rank is at most \(s\); a nonzero minor of size \(s\) forces equality. Conversely, the determinant of any nonzero minor of \(aM\) remains nonzero outside the finite set of primes dividing its principal ideal. A minor nonzero after reduction has nonzero determinant in \(K\). Multiplication by \(a\) does not change rank over \(K\).
\end{proof}

For the separating list of Proposition~\ref{prop:fitted-grid-base-separation}, the coefficient comparison gives a nonzero minor of size \(s\) on the realized base row. The determinant test also permits a particular candidate list to be checked after reduction at one finite prime. The finite and real size estimates below are taken on the same list of tests.

\subsubsection{Constituent estimates on the same matrix}

Separation identifies the row at the base tuple. To bound its height one now needs denominator and archimedean estimates for entries of this very matrix.

Write \(q_P\) for the residue cardinality at a finite place of the CM field and \(c_P\) for the coset depth there. Use
the sum \(\Sigma\) of depths times log norms of the base
moduli as in Section~\ref{subsec:fitted-row-height}.
In the local estimates below, the test and tuple coordinates are
the ones used to construct \(V(t)\).

\begin{proposition}[Finite and complex factors of the fitted matrix]\label{prop:fitted-grid-constituent-costs}
For the matrix of Proposition~\ref{prop:fitted-grid-coefficient-matrix},
at \(w\mid\ell\) use its assigned \(t(w)\) and the per-link
finite clearing of the two chains in the same coordinates.
Writing \(C_{\rm fin}\) for the weighted cost of the finite coset and moved-data factors, one has
\[
 C_{\rm fin}=O\!\left(fm^2N^2b(\log b+\Sigma)\right).
\]
 At each complex embedding of \(K\) use the base tuple \(t_0\).
The kernel and terminal comparisons use their original
coordinates in the same test. The class count is taken on the
\(2m+1\) definite vertices of the two chains, whose ranks are
at most \(N\); the selected modulus has norm at most
\(\exp(2\Sigma)\). The constants work for the whole selected grid.
\end{proposition}

\begin{proof}
\proofstep{Restrict denominator supports to the base places.}

At a finite coefficient place \(w\) above \(\ell\), the tuple
used is \(t(w)\). Its character is the deletion character times
an auxiliary character trivial at the places above \(\ell\).
Consequently a ramification place of the tuple off the base
set lies over a different rational prime. At the base tuple
this follows from the trivial twist. Outside the base set the
local coset coefficients have denominators supported at their
own residue characteristics. One way to see the support is to
filter the unipotent group by the ordered flag. For a unipotent
matrix \(g=1+U\) with \(U^d=0\), if \(g^a\) is integral and
\(a\) is prime to the residue characteristic, write \(V=g^a-1\).

Since \(V\) is a multiple of \(U\) in the
algebra generated by \(U\), one has \(V^d=0\) and
\[
        g=\sum_{j=0}^{d-1}\binom{1/a}{j}V^j .
\]
The binomial polynomial \(\binom{x}{j}\) takes integral values on
nonnegative integers and is continuous on the local field. Their
density in \(\mathbb Z_p\), for the residue characteristic \(p\),
gives \(\binom{1/a}{j}\in\mathbb Z_p\). Thus \(g\) is integral. Thus the finite relative
indices in these congruence subgroups are powers of the local
residue characteristic. Their contributions away from that
characteristic are units. At a base split place the supported
term has the staircase representative with diagonal exponents
bounded by \(mc_P\). Conjugating the standard lattice gives the
volume ratio \(q_P^z\) with \(|z|\le Nm^2c_P\) at each step:
the two lattices are compared by the sum of the absolute
diagonal orders.

Multiplying the intermediate volume factors
gives \(|z|\le Nm^3c_P\).
Writing \(z_j\) for the exponent at the \(j\)-th step gives
\[
  \left|\sum_{j=1}^{m} z_j\right|\le\sum_{j=1}^{m}|z_j|\le m\,Nm^2 c_P=Nm^3c_P .
\]
 Any root of unity in the moved
tensor datum is a unit. The nonunit factors of its Fourier
normalization come from the Jacobian of the finite polarization
changes. In the two frames the changed congruence orders are
bounded by a constant times \(m^2N^2c_P\), from the box
indices and the ranks of the matrix spaces. The costs for coefficients and data are sums of nonnegative
multiples of \(v_w(q_P)\), using the base places and the stated
depth bounds. For each such contribution the normalization gives
\[
 \sum_w\omega_w\,b_P v_w(q_P)=b_P\log q_P\qquad(b_P\ge0).
\]
Summing the base-place coefficients gives the total bound
\(O(fm^2N^2b(\log b+\Sigma))\).
\label{cor:fitted-evaluation-3}

\proofstep{Bound the volume scalar and inverse stabilizer orders.}

The common scalar is a product of powers of the square roots
of residue norms. The absolute exponents are bounded by
\(16m^3N^2\) times the sum of the absolute orders of the two dual ideals of the CM field. The weighted log sum of those
orders is bounded by the logarithm of the absolute discriminant of the CM field
on the chosen model. This controls the finite denominator
and all complex absolute values of the same scalar. For an inverse stabilizer use its faithful rational representation of
dimension at most \(D=[E:\Q]N\) and add trivial coordinates to
make the dimension \(D\).

For an initial integral lattice \(L\) and the finite stabilizer \(G\), take
\[
  L_G=\sum_{g\in G}gL,\qquad hL_G=\sum_{g\in G}hgL=L_G\quad(h\in G).
\]
 Reduction at an odd
prime \(p\) is injective on the group. Otherwise the kernel contains
an element \(1+p^kA\) of prime order \(r\), with \(k>0\) and
\(A\not\equiv0\bmod p\). The first nonzero term in its \(r\)-th
power minus one has order \(k\) for \(r\ne p\) and \(k+1\)
for \(r=p\), by binomial expansion (\(p\) is odd). Neither is
compatible with order \(r\).

For odd \(\ell\) choose an odd \(p\) primitive modulo \(\ell^2\)
by Dirichlet's theorem. Write 
\[
  h=\lfloor D/(\ell-1)\rfloor.
\]
Only \(j=(\ell-1)a\), \(1\le a\le h\), contribute factors
\(p^j-1\) to the \(\ell\)-part of \(\#\mathrm{GL}_D(\mathbb F_p)\).
Their valuations are \(1+v_\ell(a)\). Thus
\[
 v_\ell(\#\mathrm{GL}_D(\mathbb F_p))
     =h+v_\ell(h!)=v_\ell((\ell h)!)\le v_\ell((2D)!),
\]
using \(\ell h\le2D\). At \(\ell=2\) take \(p=3\) and use
\[
 v_2(3^j-1)=\begin{cases}1,&j\text{ odd},\\2+v_2(j),&j\text{ even}.\end{cases}
 \qquad \sum_{j=1}^D v_2(3^j-1)\le 2D+v_2(D!).
\]
These bounds for every \(\ell\) put each stabilizer order
among the divisors of \(8^D(2D)!\). At most \(2m+1\) inverse
stabilizers occur in a class term, so their finite mean
log cost is 
\[
  (2m+1)\log(8^D(2D)!)=O(mD\log(2D)).
\]
Their reciprocals have absolute value at most one at infinity.
Together with the common scalar estimate this is absorbed
in the allowances of Section~\ref{subsec:fitted-row-height}.

\label{cor:fitted-evaluation-4}

\proofstep{Use the same local depths, jets and terminal coefficient pairing.}

For the chains use the split and opposite-sign jet maps
in the vertex lattices. The flat product is periodic at a
depth whose order at \(\ell\) is bounded by the local base
modulus when the extra ramification there is absent. A
period at a product of local levels gives a coset expansion
at that depth. Given two such periods, the local additive
subgroups that they generate give the period at their
greatest common divisor.
For integer periods \(d_1,d_2\), the subgroup generated by both period subgroups is
\[
  d_1\mathbb Z+d_2\mathbb Z=\gcd(d_1,d_2)\mathbb Z .
\]
 The least positive depth giving a
coset expansion therefore divides every admissible depth.
This gives the bound at the depth of the expansion used
for the matrix.

The normalized finite jet estimates and the
coefficient-lattice change in Section~\ref{subsec:selected-kernel-jets}
give a single local clearing budget per link for each chain,
uniform over its occurring data and classes. The top test
is integral in these same coordinates. The initial row is
evaluated on its integral form lattice, with the inverse frame
clearing its defect. Let \(b_{\rm br}>0\) clear the branching
pairing \(B\) at the terminal node at every prime. If
\(\pi_w^{e_1}v_1,\pi_w^{e_2}v_2\) are integral in the two
local input lattices, then
\[
 v_w(B(v_1,v_2))\ge-e_1-e_2-v_w(b_{\rm br})
\]
in the discrete valuation normalized by \(v_w(\pi_w)=1\).
The common integer satisfies \(\log b_{\rm br}=O(fN^2 b\log b)\)
under the interlacing weights. This gives the single terminal
charge alongside the per-link costs. \label{cor:fitted-evaluation-5}

\proofstep{Count the support and bound the matrix at the base tuple.}

At infinity use the base tuple. The local absolute bounds on
the scalar, volume and Fourier factors use the same depths.
The maximum support count is obtained from the finite
quotients defining the data and the two coset decompositions.
The class counts use the product of the definite class
quotients on the \(2m+1\) vertices, with their modulus indices.
For one vertex write \(h_1\) for the number of classes at level one. In rank at most \(N\) the class estimate gives
\[
  h_1\le(12f^2N^4\operatorname{rd}(E)^2)^{6fN^2}.
\]
The congruence quotient injects into the matrix reduction at
the selected modulus. Its cardinality is at most
\(\exp(2\Sigma N^2)\), since the logarithm of the modulus
norm is at most \(2\Sigma\).

The size bounds give \(N\le b\) and
\(\log(12f^2N^4\operatorname{rd}(E)^2)\le8\log b\).
The logarithm of the product over the \(2m+1\)
vertices is therefore at most
\[
 (2m+1)(48fN^2\log b+2N^2\Sigma)
           \le m(144fNb\log b+6N^2\Sigma),
\]
using \(2m+1\le3m\) and \(N\le b\).
The initial test-class count uses the single-vertex bound.
The datum and coset counts and these class counts contribute
the three sum bounds and the test block bound used below.
The chain operator estimates use the
Hermitian norm of the CM-normalized kernels at each
embedding, including the opposite conjugate reading for the
row chain. Their log gains can depend on the embedding.

At a fixed embedding of \(K\), the algebraic matrix of each
kernel is fixed. Its bound from the period and theta-norm
comparison holds for every extension of the embedding to the
curve coefficient field. Averaging the log bounds over those
extensions keeps a bound for the same matrix.
If \(g_{\tilde v}\) bounds the log norm at each extension of a fixed embedding \(v\), averaging preserves the inequality
\[
  \log\|A_v\|\le \frac{1}{\#\{\tilde v\mid v\}}\sum_{\tilde v\mid v}g_{\tilde v}.
\]
 The resulting
mean is bounded using the CM area and theta norm on that
curve model. The terminal branching estimate is a squared
operator bound, and the initial form is measured against
its actual row Gram. This gives the constituent estimates
with a common comparison constant at \(t_0\) and no need
for a bound on every complex embedding of a twisted tuple.
\label{cor:fitted-evaluation-6}
\end{proof}
